\documentclass{SGGW-thesis}

\MAGISTERSKAtrue
\WZIMtrue

% ============================================================================
% METADANE PRACY — uzupełnij przed pierwszą kompilacją
% ============================================================================

\title{Analiza podstawowych współczynników biegu\\przy pomocy uczenia maszynowego}
\Etitle{Analysis of basic running coefficients\\using machine learning}

\author{Imię Nazwisko}                     % TODO: uzupełnij
\date{2026}                                % TODO: uzupełnij rok obrony
\album{000000}                             % TODO: numer albumu
\thesis{Praca dyplomowa magisterska na kierunku:}
\course{Informatyka}                       % TODO: potwierdź kierunek
\promotor{<tytuł> <Imię Nazwisko> promotora} % TODO: dane promotora
\pworkplace{Instytut Informatyki Technicznej\\Katedra <nazwa katedry>} % TODO

% ============================================================================
% PAKIETY DODATKOWE
% ============================================================================

\usepackage{hyperref}
\usepackage{graphicx}
\usepackage{booktabs}      % ładne tabele (\toprule, \midrule, \bottomrule)
\usepackage{amsmath}       % wzory matematyczne
\usepackage{listings}      % bloki kodu
\usepackage{xcolor}        % kolory (np. dla kodu)

\begin{document}

\maketitle
\statementpage

% ============================================================================
% STRESZCZENIE PL + EN
% ============================================================================

\abstractpage
{Analiza podstawowych współczynników biegu przy pomocy uczenia maszynowego}
{
% TODO: streszczenie po polsku (~200 słów). Pisać NA KOŃCU, gdy reszta pracy gotowa.
% Powinno zawierać: cel pracy, metoda (MediaPipe + LSTM + reguły),
% główne wyniki (accuracy 70.9\%, 12 współczynników, 13+ reguł), wkład.
Streszczenie pracy w języku polskim.
}
{uczenie maszynowe, biomechanika biegu, MediaPipe, klasyfikacja faz biegu, LSTM, monocular 2D wideo}
{Analysis of basic running coefficients using machine learning}
{
% TODO: abstract in English (~200 words)
Thesis abstract in English.
}
{machine learning, running biomechanics, MediaPipe, running phase classification, LSTM, monocular 2D video}

% ============================================================================
% SPIS TREŚCI
% ============================================================================

{
  \doublespacing
  \tableofcontents
}

\startchapterfromoddpage

% ============================================================================
% ROZDZIAŁY (każdy w osobnym pliku w chapters/)
% ============================================================================

\chapter{Wstęp}
\label{ch:wstep}

\section{Motywacja}
\label{sec:motywacja}

Bieganie należy do najpopularniejszych form aktywności fizycznej, a~wraz z~jego rosnącą popularnością rośnie liczba kontuzji przeciążeniowych związanych z~techniką biegu. Część tych urazów wiąże się z~parametrami, które można zmierzyć i~skorygować --- zbyt niską kadencją, nadmiernym wykrokiem (overstridingiem) czy asymetrią obciążenia nóg. Obiektywna analiza techniki biegu mogłaby zatem wspierać prewencję kontuzji, pod warunkiem że jest dostępna poza wyspecjalizowanym laboratorium.

Tradycyjna analiza biomechaniczna opiera się na systemach motion capture z~markerami, które oferują wysoką dokładność, lecz są kosztowne, wymagają laboratorium i~obecności specjalisty --- przez co pozostają poza zasięgiem przeciętnego biegacza. Jednocześnie sami biegacze słabo oceniają własną technikę: badanie przeprowadzone na grupie 710~biegaczy wykazało, że trafność samodzielnego rozpoznania własnego wzorca lądowania stopy wynosi zaledwie około 43\%~\cite{vincent2024footstrike}. Wskazuje to na potrzebę narzędzia dostarczającego obiektywnej oceny techniki w~sposób tani i~dostępny.

Rozwój metod automatycznego wykrywania ułożenia ciała na nagraniu wideo --- bez przyklejania do niego jakichkolwiek fizycznych markerów otworzył taką możliwość. Z~nagrania ze~zwykłej kamery można dziś wykryć położenie punktów charakterystycznych ciała i~na tej podstawie analizować ruch. Niniejsza praca wykorzystuje to podejście do zbudowania systemu, który z~pojedynczego nagrania biegu na bieżni mechanicznej wylicza współczynniki biomechaniczne i~generuje rekomendacje treningowe. 
\newpage
\section{Problem badawczy}
\label{sec:problem-badawczy}

Praca odpowiada na trzy główne pytania badawcze.

Pierwsze pytanie dotyczy projektu i~dokładności samego systemu, jak zaprojektować pipeline uczenia maszynowego --- prowadzący od estymacji pozy, przez klasyfikację faz biegu, po obliczenia geometryczne --- który z~pojedynczego nagrania jedną kamerą wylicza komplet współczynników biegu, i~jaką dokładność osiąga dla poszczególnych z~nich.

Drugie pytanie dotyczy klasyfikatora faz biegu jako fundamentu całego pipeline'u. Który z~rozważanych klasyfikatorów najtrafniej rozpoznaje fazy kontaktu z~podłożem oraz fazy lotu, a~także jak dokładność tego rozpoznawania przekłada się na precyzję wyliczanych z~faz współczynników czasowych, takich jak kadencja czy czas kontaktu z~podłożem.

Trzecie pytanie dotyczy tego, jak warunki nagrania wpływają na wyniki. Które współczynniki są wrażliwe na to, czy nagranie jest poziome czy pionowe, na kąt ustawienia kamery oraz na liczbę klatek na sekundę, a~także czy można automatycznie rozpoznać przypadki, w~których dany współczynnik staje się niewiarygodny z~powodu sposobu nagrania.

Przy projektowaniu pracy przyjąłem kilka założeń, które następnie zweryfikowałem. Spodziewałem się, że model sekwencyjny, uwzględniający kontekst sąsiednich klatek, przewyższy klasyfikator analizujący każdą klatkę niezależnie, i~że dokładniejsza klasyfikacja przełoży się na precyzyjniejsze metryki czasowe. Zakładałem również, że współczynniki przestrzenne, oparte na geometrii ciała w~płaszczyźnie obrazu, okażą się wrażliwe na perspektywę kamery, podczas gdy współczynniki czasowe, oparte na czasie trwania faz, będą wrażliwe przede wszystkim na częstotliwość klatek. Przyjąłem wreszcie, że najbardziej zawodne przypadki uda się rozpoznać automatycznie, bez dodatkowego pomiaru wzorcowego, sprawdzając, czy wyliczona wartość mieści się w~zakresie możliwym dla człowieka. 

\section{Cel i~zakres pracy}
\label{sec:cel-zakres}

Celem pracy jest zaprojektowanie, implementacja i~ocena systemu, który z~nagrania biegu na bieżni oblicza współczynniki biomechaniczne i~generuje na ich podstawie rekomendacje treningowe oparte na literaturze.
\newpage
Pierwotnie zakładałem analizę kilku podstawowych współczynników --- kadencji, czasu kontaktu z~podłożem, długości kroku i~wzorca lądowania stopy. W~trakcie realizacji okazało się, że z~tych samych danych wejściowych (sekwencji faz i~wygładzonych keypointów) da się w~naturalny sposób obliczyć więcej wielkości, bez dodatkowych wymagań wobec nagrania. Ostateczny zakres objął dwanaście współczynników --- pięć czasowych i~siedem przestrzennych --- uzupełnionych o~wskaźniki symetrii między lewą a~prawą stroną ciała.

Uczenie maszynowe wykorzystuję w~jednym, ale kluczowym miejscu pipeline'u --- na etapie klasyfikacji faz biegu. To wytrenowany model rozstrzyga dla każdej klatki, czy biegacz znajduje się w~fazie kontaktu lewej stopy, prawej stopy czy lotu. Bez tej klasyfikacji nie da się wyznaczyć żadnego ze~współczynników czasowych, więc jest ona warunkiem działania całego systemu. Same współczynniki obliczam następnie metodami geometrycznymi z~położenia keypointów, a~rekomendacje generuję regułami opartymi na literaturze biomechanicznej, a~nie modelem uczonym z~danych. Uczenie maszynowe jest więc narzędziem umożliwiającym całą analizę, choć nie jedynym jej elementem.

\chapter{Stan wiedzy}
\label{ch:sota}

Rozdział ten przedstawia tło literaturowe pracy w~czterech obszarach: estymacji pozy z~wideo monocular 2D, biomechaniki biegu, klasyfikacji faz biegu i~chodu metodami uczenia maszynowego oraz walidacji metod 2D względem trójwymiarowego motion capture. Na tej podstawie identyfikuję lukę badawczą, którą praca wypełnia. Szczegóły metod zastosowanych w~niniejszej pracy opisuję w~rozdziale~\ref{ch:materialy-metody}; tutaj koncentruję się na dorobku innych autorów.

\section{Estymacja pozy z~wideo monocular 2D}
\label{sec:sota-pose}

Klasyczna analiza ruchu opiera się na systemach motion capture z~markerami (np.~Vicon, Qualisys), które oferują wysoką dokładność, lecz są drogie, wymagają laboratorium i~doświadczonego operatora. Rozwój metod estymacji pozy bez markerów (\textit{markerless pose estimation}) otworzył możliwość analizy ruchu z~nagrań ze~zwykłej kamery. Frameworki takie jak OpenPose~\cite{cao2021openpose} oraz MediaPipe BlazePose~\cite{bazarevsky2020blazepose}, trenowane m.in.\ na zbiorze COCO~\cite{lin2014coco}, wykrywają keypointy ciała w~każdej klatce wideo, działając na pojedynczej kamerze (\textit{monocular}) i~obrazie dwuwymiarowym.

Roggio i~wsp.~\cite{roggio2024review} w~przeglądzie dziewięciu modeli estymacji pozy wskazują na ich potencjał w~zastosowaniach klinicznych, sportowych i~ergonomicznych jako tania, nieinwazyjna alternatywa dla systemów markerowych, zaznaczając jednocześnie wyzwania związane z~dokładnością, jakością danych i~zasłonięciami. Edriss i~wsp.~\cite{frontiers2025review} w~przeglądzie poświęconym analizie ruchu w~sporcie podkreślają, że rozwiązania 2D są ekonomiczne i~proste w~użyciu, lecz wciąż mają trudność z~rejestracją ruchów poza płaszczyzną obrazu (\textit{out-of-plane}). Porównanie konkretnych standardów keypointów oraz uzasadnienie wyboru MediaPipe Pose w~niniejszej pracy przedstawiam w~sekcji~\ref{sec:mediapipe}.

\section{Biomechanika biegu --- współczynniki i~ich znaczenie}
\label{sec:sota-wspolczynniki}

Cykl biegowy różni się od cyklu chodu: w~biegu występuje faza lotu (obie stopy w~powietrzu), a~zanika faza podwójnego podparcia charakterystyczna dla chodu~\cite{novacheck1998}. Z~tego powodu progi referencyjne i~wzorce ruchu ustalone dla chodu klinicznego nie przenoszą się wprost na bieg sportowy.

Literatura biomechaniczna dostarcza ugruntowanych współczynników opisujących technikę biegu. Heiderscheit i~wsp.~\cite{heiderscheit2011} wykazali, że manipulacja kadencją wpływa na obciążenia stawów i~wiąże się z~overstridingiem oraz kątem kolana w~momencie kontaktu. Souza~\cite{souza2016} omawia czas kontaktu z~podłożem i~wzorzec lądowania stopy w~kontekście analizy wideo. Novacheck~\cite{novacheck1998} opisuje fazy cyklu biegowego, pochylenie tułowia, kąty stawów i~oscylację pionową. Daoud i~wsp.~\cite{daoud2012} powiązali wzorzec lądowania stopy z~częstością kontuzji u~biegaczy długodystansowych, a~Robinson i~wsp.~\cite{robinson1987} zaproponowali wskaźnik symetrii (\textit{Symmetry Index}) do ilościowego opisu różnic między stronami ciała. Przeglądy nowsze, jak Figueiredo i~wsp.~\cite{figueiredo2025cadence}, potwierdzają związek kadencji z~prewencją kontuzji. Konkretne zakresy referencyjne tych współczynników i~sposób ich obliczania w~niniejszej pracy zestawiam w~sekcji~\ref{sec:obliczanie-wspolczynnikow}.

\section{Klasyfikacja faz biegu i~chodu metodami uczenia maszynowego}
\label{sec:sota-klasyfikatory}

Rozpoznawanie faz cyklu ruchu jest przedmiotem wielu prac, najczęściej opartych na czujnikach inercyjnych (IMU) noszonych na ciele. Vu i~wsp.~\cite{vu2020review} dokonali przeglądu algorytmów detekcji faz chodu dla protez kończyn dolnych. Su i~wsp.~\cite{su2020dcnn} zastosowali głęboką sieć konwolucyjną z~danych IMU, a~Jeon i~Lee~\cite{jeon2024bilstm} dwukierunkową sieć LSTM (BiLSTM) odporną na zmiany kierunku ruchu --- architekturę pokrewną zastosowanej w~niniejszej pracy.

W~kontekście biegu klasyczne metody uczenia maszynowego wykorzystywano m.in.\ do wykrywania zmian zmęczeniowych: Dimmick i~wsp.~\cite{dimmick2023fatigue} zastosowali las losowy (Random Forest) do analizy biomechaniki biegu, a~Martínez-Gramage i~wsp.~\cite{martinez2020rf} --- do redukcji ryzyka kontuzji u~młodych triathlonistów. Prace łączące estymację pozy z~wideo z~klasyfikacją uczenia maszynowego pojawiają się rzadziej: Ali i~wsp.~\cite{ali2024dmd} połączyli estymację pozy (m.in.\ MediaPipe) z~klasyfikatorem SVM oraz LSTM-CNN do rozróżniania patologii chodu, a~Young i~wsp.~\cite{young2023smartphone} zaproponowali ocenę biegu z~pojedynczej kamery smartfona. Większość prac dotyczy jednak chodu klinicznego i~opiera się na czujnikach noszonych, nie na czystym wideo.

\section{Walidacja estymacji 2D względem 3D motion capture}
\label{sec:sota-walidacja}

Kluczowym wątkiem literatury jest pytanie, na ile estymacja pozy z~wideo 2D dorównuje referencyjnym systemom 3D. Stenum i~wsp.~\cite{stenum2021} porównali OpenPose z~systemem 3D MoCap dla chodu, uzyskując błąd średni rzędu 0{,}02~s dla metryk czasowych oraz kilku stopni dla kątów stawów. Wskazali przy tym istotne ograniczenie: długość kroku jest estymowana najdokładniej, gdy osoba znajduje się w~centrum pola widzenia kamery --- a~więc wynik zależy od pozycji w~kadrze. Menychtas i~wsp.~\cite{menychtas2023} porównali OpenPose i~MediaPipe z~systemem Vicon, raportując błędy w~ruchach prostopadłych do płaszczyzny obrazu oraz problemy z~lokalizacją keypointu kostki. Praca oparta na metodzie HGcnMLP~\cite{hgcnmlp2024} pokazała wysoką zgodność estymacji 3D ze~smartfona z~systemem Vicon, lecz dane zbierano wyłącznie przy jednej orientacji kamery ($90^\circ$), bez walidacji innych perspektyw.

Z~prac tych wyłania się spójny obraz ograniczeń estymacji 2D: błędy dla ruchów poza płaszczyzną obrazu, zależność od perspektywy oraz od pozycji w~polu widzenia. Ograniczenia te są w~literaturze \textit{identyfikowane}, ale rzadko \textit{systematycznie analizowane} pod kątem wpływu na konkretne współczynniki --- co stanowi część luki opisanej poniżej.

\section{Luka badawcza}
\label{sec:sota-luka}

Z~przeglądu literatury wynika, że żadna z~analizowanych prac nie łączy jednocześnie następujących elementów, które wyznaczają zakres niniejszej pracy:

\begin{enumerate}
    \item \textbf{Ukierunkowanie na bieg, nie chód.} Analizowane prace dotyczą niemal wyłącznie chodu (gait), zwykle w~kontekście klinicznym. Bieg sportowy na bieżni ma odmienne fazy (występuje faza lotu), inne progi kadencji i~czasu kontaktu oraz inne wartości referencyjne --- stanowi więc odrębną niszę.
    \item \textbf{Systematyczne porównanie czterech klasyfikatorów na małym zbiorze.} Większość prac proponuje jeden model. W~niniejszej pracy porównuję cztery konfiguracje (las losowy na cechach surowych i~inżynierowanych oraz dwie sieci BiLSTM), traktując mały, dobrze opisany zbiór jako świadomą specyfikację, nie wadę.
    \item \textbf{Korekta proporcji obrazu jako krok preprocessingu.} Menychtas i~wsp.~\cite{menychtas2023} identyfikują błędy dla ruchów poza płaszczyzną obrazu, lecz nie raportują badania wpływu korekty proporcji kadru na jakość klasyfikatora faz.
    \item \textbf{Silnik reguł rekomendacji oparty na literaturze.} Prace cytują wartości referencyjne (Heiderscheit, Souza, Novacheck, Robinson, Daoud), ale nie operacjonalizują ich w~postaci wykonywalnych reguł z~poziomami ważności, źródłami i~regułami łączonymi.
    \item \textbf{Automatyczna detekcja niskiej wiarygodności pomiaru.} Stenum~\cite{stenum2021} i~Menychtas~\cite{menychtas2023} identyfikują problem perspektywy, lecz żadna z~prac nie sygnalizuje automatycznie przypadków, w~których dany współczynnik staje się nieinterpretowalny.
\end{enumerate}

Pełne sformułowanie pytań badawczych i~hipotez, wynikających z~tej luki, przedstawiam w~rozdziale~\ref{ch:wstep}.

\chapter{Materiały i metody}
\label{ch:materialy-metody}

\section{Dataset}
\label{sec:dataset}

Zbiór danych wykorzystany w~niniejszej pracy składa się z~15 nagrań wideo przedstawiających bieg na bieżni mechanicznej w~ujęciu z~boku. Materiał pochodzi z~dwóch źródeł: publicznie dostępnych nagrań opublikowanych w~serwisie YouTube oraz nagrań własnych zarejestrowanych na potrzeby projektu z~udziałem trzech członków klubu biegowego \textit{Zabiegani na Szamę} --- Pawła, Adama i~Janka --- którym dziękuje za pomoc i udział w~badaniu. Nagrania Pawła (film~11) i~Adama (film~12) włączyłem do zbioru treningowego, natomiast materiał Janka posłużył wyłącznie do walidacji pełnego pipeline'u na wytrenowanym modelu. Te 15 nagrań pochodzi od~kilkunastu biegaczy o~zróżnicowanym poziomie zaawansowania --- od rekreacyjnych joggerów po zawodowego ultramaratończyka.

\subsection{Kryteria doboru nagrań}
\label{sec:dataset-kryteria}

Minimalne kryteria włączenia nagrania do zbioru były następujące:

\begin{itemize}
    \item ujęcie z~boku  --- kamera ustawiona w~przybliżeniu prostopadle do~kierunku biegu,
    \item bieżnia mechaniczna --- stałe tło umożliwiające stabilną detekcję pozy,
    \item sylwetka biegacza widoczna w~kadrze w~stopniu wystarczającym do detekcji keypointów,
    \item co najmniej 10~sekund ciągłego biegu bez zasłonięcia sylwetki,
    \item jakość obrazu wystarczająca do poprawnej detekcji keypointów przez MediaPipe Pose.
\end{itemize}

Do zbioru celowo włączyłem nagrania odbiegające od warunków idealnych --- m.in.\ wideo w~orientacji pionowej (film~10), nagrania w~trybie slow-motion o~efektywnym FPS poniżej 15 (filmy~4, 7, 8) --- w~celu oceny odporności pipeline'u na zróżnicowane warunki nagrania.

\subsection{Charakterystyka zbioru}
\label{sec:dataset-charakterystyka}

Tabela~\ref{tab:dataset} przedstawia parametry techniczne wszystkich nagrań w~zbiorze danych. Materiał charakteryzuje się znacznym zróżnicowaniem pod względem rozdzielczości (od $320 \times 360$ do $1920 \times 1080$~pikseli), liczby klatek na sekundę (od~9{,}46 do~30~) oraz proporcji obrazu (ujęcia poziome i~pionowe). Część nagrań zarejestrowana została w~trybie slow-motion (filmy~4, 7, 8, 13), co skutkuje niższym efektywnym FPS w~metadanych pliku wideo.

\begin{table}[htbp]
\centering
\caption{Parametry techniczne nagrań w~zbiorze danych}
\label{tab:dataset}
\small
\begin{tabular}{clrrrl}
\hline
\textbf{Nr} & \textbf{Opis} & \textbf{FPS} & \textbf{Rozdzielczość} & \textbf{Klatki} & \textbf{Uwagi} \\
\hline
1  & Bieg w~obuwiu              & 29{,}97 & $480 \times 360$   & 949  & --- \\
2  & Bieg 13\,km/h              & 30{,}00 & $360 \times 450$   & 300  & --- \\
3  & Bieg 15\,km/h              & 30{,}00 & $360 \times 450$   & 300  & --- \\
4  & Bieg 15\,km/h              &  9{,}46 & $360 \times 450$   & 560  & slow-motion \\
5a & Bieg rekreacyjny (seg.\ 1) & 29{,}97 & $320 \times 360$   & 660  & --- \\
5b & Bieg rekreacyjny (seg.\ 2) & 29{,}97 & $320 \times 360$   & 300  & --- \\
6a & Bieg maratoński (seg.\ 1)  & 29{,}97 & $512 \times 360$   & 892  & --- \\
6b & Bieg maratoński (seg.\ 2)  & 29{,}97 & $512 \times 360$   & 438  & --- \\
7  & Bieg 4\,m/s                & 13{,}33 & $360 \times 450$   & 800  & slow-motion \\
8  & Bieg 16\,km/h         & 11{,}25 & $364 \times 360$   & 900  & slow-motion \\
9  & Bieg zmiennym tempem       & 30{,}00 & $640 \times 360$   & 870  & --- \\
10 & Bieg kamera pionowa       & 29{,}97 & $608 \times 1080$  & 320  & wideo pionowe \\
11 & Paweł --- bieg             & 29{,}97 & $1280 \times 720$  & 1649 & nagranie własne \\
12 & Adam --- bieg              & 29{,}99 & $1920 \times 1080$ & 2399 & nagranie własne \\
13 & Bieg 2 4\,m/s                 & 15{,}00 & $600 \times 800$   & 750  & slow-motion \\
\hline
\multicolumn{4}{l}{\textit{Razem:}} & \textit{12\,087} & \\
\hline
\end{tabular}
\end{table}

Cztery nagrania w~zbiorze (filmy~4, 7, 8, 13) zarejestrowane zostały w~trybie slow-motion, co przekłada się na niski FPS w~metadanych pliku wideo (od~9{,}46 do~15~FPS). Pomimo pozornie niskiej płynności, nagrania te są pełnowartościowym materiałem treningowym --- ruch biegacza jest rozciągnięty w~czasie, dzięki czemu jeden cykl biegowy obejmuje kilkadziesiąt klatek zamiast kilkunastu. Klasyfikator uczy się rozpoznawać pozycję ciała na podstawie geometrii keypointów, nie na podstawie czasu między klatkami, dlatego większa ilość klatek na cykl stanowi dodatkowy atut.

Łącznie 15~nagrań obejmuje $\approx$~12\,000~klatek, z~których większość posiada etykiety faz biegu. Zbiór podzieliłem na część treningową (10~filmów, $\approx$~9\,800~klatek), walidacyjną (2~filmy, $\approx$~740~klatek) i~testową (3~filmy, $\approx$~1\,500~klatek).

\section{Pipeline ekstrakcji keypointów}
\label{sec:pipeline-extraction}

Ekstrakcja punktów charakterystycznych ciała (keypointów) z~nagrań wideo stanowi pierwszy etap pipeline'u analitycznego. Proces przebiega w~trzech krokach: dekodowanie klatek wideo (biblioteka OpenCV~\cite{opencv}), detekcja pozy przez sieć neuronową oraz wygładzanie uzyskanych trajektorii keypointów.

\subsection{Detekcja pozy --- MediaPipe Pose}
\label{sec:mediapipe}

Do detekcji pozy zastosowałem bibliotekę MediaPipe Pose w~wersji 0.10.14 (legacy Solutions API) \cite{bazarevsky2020blazepose}. Model wykrywa 33~keypointy w~każdej klatce wideo: 1~punkt centralny (nos) oraz 16~par punktów symetrycznych L/P obejmujących twarz, tułów, kończyny górne i~dolne --- w~tym pięty (\texttt{HEEL}) i~czubki stóp (\texttt{FOOT\_INDEX}). Każdy keypoint opisany jest czterema wartościami:

\begin{itemize}
    \item $x, y \in [0, 1]$ --- znormalizowane współrzędne w~płaszczyźnie obrazu (0{,}0 = lewy górny róg, 1{,}1 = prawy dolny róg); oś $x$ rośnie w~prawo, a~oś $y$ w~dół,
    \item $z$ --- estymowana współrzędna głębi, określająca przybliżone położenie keypointu względem kamery. Ponieważ nagrania pochodzą z pojedynczej kamery, wartość $z$ nie jest wynikiem bezpośredniego pomiaru odległości, lecz predykcją modelu MediaPipe. Z tego powodu traktowałem ją jako informację pomocniczą, mniej stabilną i trudniejszą do interpretacji niż współrzędne $x$ i $y$. W dalszych etapach pracy wykorzystywałem głównie dane dwuwymiarowe.

    \item $visibility \in [0, 1]$ --- wskaźnik pewności detekcji, określający prawdopodobieństwo, że dany keypoint jest widoczny w~kadrze i~nie jest zasłonięty przez inne części ciała. Sieć neuronowa uczy się przewidywać tę wartość na podstawie tego, czy okolica keypointu jest przesłonięta przez inne kończyny lub wykracza poza kadr, oraz czy cechy wizualne w~otoczeniu keypointu są spójne z~oczekiwaną konfiguracją pozy. Wartość generowana jest przez warstwę sigmoidalną na wyjściu modelu.
\end{itemize}

Zastosowałem najwyższy poziom złożoności modelu (\texttt{model\_complexity=2}), który zapewnia najdokładniejszą lokalizację keypointów kosztem dłuższego czasu przetwarzania (inferencji). Tryb wideo (\texttt{static\_image\_mode=False}) umożliwia wykorzystanie śledzenia między klatkami, co poprawia stabilność trajektorii. Przykład pozy wykrytej tą metodą na klatce nagrania przedstawia rysunek~\ref{fig:szkielet-mediapipe}.

\begin{figure}[H]
\centering
\includegraphics[width=0.45\textwidth]{figures/szkielet_13kmh_prawa.png}
\caption{Poza wykryta przez MediaPipe Pose --- 33~keypointy połączone w~model szkieletowy, nałożone na klatkę nagrania biegu w~ujęciu bocznym.}
\label{fig:szkielet-mediapipe}
\end{figure}

MediaPipe Pose wybrałem spośród trzech rozważanych standardów detekcji pozy:

\begin{itemize}
    \item \textbf{COCO} \cite{lin2014coco} --- benchmark referencyjny definiujący 17~keypointów ciała (nos, oczy, uszy, ramiona, łokcie, nadgarstki, biodra, kolana, kostki). Nie zawiera żadnych keypointów stopy poniżej kostki,
    \item \textbf{OpenPose BODY\_25} \cite{cao2021openpose} --- rozszerzenie standardu COCO do 25~keypointów, dodające m.in.\ pięty (\texttt{HEEL}), paluch (\texttt{BIG\_TOE}) i~mały palec (\texttt{SMALL\_TOE}) dla każdej stopy,
    \item \textbf{MediaPipe BlazePose} \cite{bazarevsky2020blazepose} --- 33~keypointy, w~tym pięty (\texttt{HEEL}) oraz \texttt{FOOT\_INDEX} (punkt odpowiadający przedniej części stopy).
\end{itemize}

Choć OpenPose również oferuje keypointy stopy, wybrałem MediaPipe z~dwóch powodów. Po~pierwsze, punkt \texttt{FOOT\_INDEX} lepiej reprezentuje przednią krawędź kontaktu stopy z~podłożem niż \texttt{BIG\_TOE}, co jest istotne przy wyznaczaniu kąta lądowania stopy. Po~drugie, MediaPipe jest zoptymalizowany pod detekcję pojedynczej osoby i~nie wymaga karty graficznej --- wystarczy standardowy procesor, co upraszcza wdrożenie pipeline'u.

Z~33 wykrywanych keypointów w~dalszej analizie wykorzystywałem przede wszystkim 13~punktów istotnych dla biomechaniki biegu: nos, ramiona (2), biodra (2), kolana (2), kostki (2), pięty (2) i~czubki stóp (2). Pozostałe keypointy (twarz, dłonie, łokcie, nadgarstki) rejestrowałem, lecz nie wnosiły informacji do klasyfikacji faz ani obliczania współczynników.

\subsection{Wygładzanie trajektorii --- filtr Savitzky-Golay}
\label{sec:savgol}

\subsubsection{Problem szumu w~trajektoriach keypointów}

Każda współrzędna keypointu wyznaczonego przez MediaPipe Pose tworzy ciąg wartości zmieniających się w~czasie, czyli trajektorię opisującą ruch danego punktu ciała, np.~kolana lub pięty, z~klatki na klatkę. MediaPipe przetwarza każdą klatkę wideo w~znacznej mierze niezależnie: chociaż tryb wideo wykorzystuje mechanizm śledzenia pozy między kolejnymi klatkami, pozycja każdego keypointu jest ostatecznie wynikiem predykcji sieci neuronowej dla danej klatki. Ponieważ predykcja ta nie jest idealnie powtarzalna, pozycja keypointu może oscylować o~kilka pikseli między kolejnymi klatkami nawet wtedy, gdy rzeczywisty ruch ciała jest płynny. Zjawisko to określam jako szum lub \textit{jitter}, czyli drobne, przypadkowe odchylenia wykrytej pozycji keypointu od jego rzeczywistego położenia.

Obliczanie współczynników biomechanicznych, takich jak kąty stawów czy czas kontaktu stopy, bezpośrednio na surowych, zaszumionych trajektoriach prowadziłoby do błędów wynikających z~obecności jittera. Konieczne jest zatem wygładzenie sygnału, czyli ograniczenie losowych fluktuacji przy jednoczesnym zachowaniu rzeczywistego przebiegu ruchu \cite{crenna2021filtering}.

\subsubsection{Zasada działania filtru Savitzky-Golay}

Do wygładzania zastosowałem filtr Savitzky-Golay \cite{savitzky1964}. Dla każdego punktu w~trajektorii filtr bierze okno $w$ sąsiednich próbek, np.~$w = 11$ kolejnych klatek, i~dopasowuje do nich wielomian o~zadanym rzędzie $p$ metodą najmniejszych kwadratów. Wartość wygładzona odpowiada wartości tego wielomianu w~środkowym punkcie okna. Następnie okno przesuwa się o~jedną klatkę i~procedura powtarza się dla kolejnej próbki sygnału.

Zachowanie filtru kontrolują dwa główne parametry:

\begin{itemize}
    \item \textbf{Długość okna} $w = 11$ --- liczba kolejnych klatek branych pod uwagę przy każdym dopasowaniu. Większe okno powoduje silniejsze wygładzenie sygnału, ponieważ więcej próbek uczestniczy w~lokalnym dopasowaniu wielomianu, ale jednocześnie zwiększa ryzyko rozmycia szybkich zmian trajektorii. Przy 30~FPS okno 11~klatek obejmuje około 0{,}37~s sygnału, przy czym względem analizowanej klatki uwzględnianych jest po pięć klatek wstecz i~pięć klatek w~przód, co odpowiada około 0{,}17~s z~każdej strony. Parametr ten dobrałem jako kompromis między redukcją jittera a~zachowaniem lokalnych ekstremów istotnych dla detekcji faz biegu.
    
    \item \textbf{Rząd wielomianu} $p = 3$ --- stopień wielomianu dopasowywanego w~każdym oknie. Wielomian 3.~stopnia potrafi odwzorować lokalną krzywiznę sygnału, w~tym jego lokalne maksima i~minima. Jest to przewaga nad prostą średnią ruchomą, która ma tendencję do spłaszczania ekstremów. Zachowanie lokalnych ekstremów jest istotne, ponieważ momenty kontaktu stopy z~podłożem mogą objawiać się jako charakterystyczne maksima lub minima w~trajektoriach keypointów stopy.
\end{itemize}

Wygładzanie stosowałem oddzielnie dla współrzędnych $x$, $y$ i~$z$ każdego z~33~keypointów. Wartości $visibility$ nie wygładzałem, ponieważ jest to miara pewności detekcji generowana przez sieć neuronową, a~nie sygnał opisujący fizyczny ruch ciała.

\subsubsection{Uzupełnianie brakujących detekcji}

W~pojedynczych klatkach MediaPipe może nie wykryć pozy, np.~z~powodu rozmycia ruchu, chwilowego zasłonięcia części ciała lub niejednoznacznego ułożenia sylwetki. Filtr Savitzky-Golay wymaga ciągłego sygnału, dlatego przed jego zastosowaniem konieczne było uzupełnienie brakujących wartości.

Luki w~trajektoriach uzupełniałem interpolacją liniową. Jeżeli keypoint miał znaną pozycję w~klatce $a$ oraz w~klatce $b$, a~brakowało danych w~klatkach pomiędzy nimi, to brakujące wartości wyznaczałem jako punkty leżące na odcinku łączącym pozycję z~klatki $a$ z~pozycją z~klatki $b$. Interpolację stosowałem dla luk wewnętrznych, czyli takich, dla których istniała znana pozycja keypointu zarówno przed, jak i~po brakującym fragmencie.

\subsubsection{Ocena skuteczności wygładzania}

Skuteczność filtracji oceniałem za pomocą metryki \textit{jittera}, definiowanej jako odchylenie standardowe drugiej różnicy sygnału. Druga różnica ciągu $x_0, x_1, \ldots, x_n$ dana jest wzorem:
\[
d^2_i = x_{i+2} - 2 \cdot x_{i+1} + x_i
\]
i~stanowi dyskretne przybliżenie lokalnej zmiany prędkości sygnału, czyli jego przyspieszenia. Dla gładkiej trajektorii wartości drugiej różnicy zmieniają się w~sposób uporządkowany, dlatego ich odchylenie standardowe jest relatywnie małe. \\ Dla sygnału zaszumionego druga różnica przyjmuje bardziej nieregularne wartości, co prowadzi do wzrostu $\mathrm{std}(d^2)$.

Metrykę tę obliczałem dla współrzędnych $x$ i~$y$ dwunastu kluczowych keypointów: ramion, bioder, kolan, kostek, pięt oraz czubków stóp. Porównanie wartości metryki przed i~po filtracji pozwalało ilościowo ocenić, w~jakim stopniu filtr Savitzky-Golay ograniczył wysokoczęstotliwościowe fluktuacje trajektorii. Procentową redukcję jittera wyznaczałem ze wzoru:
\[
\text{redukcja} =
\frac{
\text{jitter}_{\text{raw}} - \text{jitter}_{\text{smooth}}
}{
\text{jitter}_{\text{raw}}
}
\cdot 100\%.
\]
Metryka ta stanowi prostą miarę poziomu wysokoczęstotliwościowych fluktuacji sygnału i~pozwala porównać trajektorie przed oraz po zastosowaniu filtracji.


\subsection{Format wyjściowy}
\label{sec:extraction-output}

Wynikiem ekstrakcji jest jeden plik CSV na nagranie, zawierający kolumny: numer klatki, znacznik czasowy (obliczony z~FPS), flaga detekcji pozy oraz $33 \times 4 = 132$ kolumny z~wygładzonymi keypointami (\texttt{\{nazwa\}\_x}, \texttt{\{nazwa\}\_y}, \texttt{\{nazwa\}\_z}, \texttt{\{nazwa\}\_visibility}).

\section{Auto-etykietowanie faz biegu}
\label{sec:auto-labeling}

Klasyfikator faz biegu wymaga etykiet treningowych --- dla każdej klatki wideo informacji, czy biegacz znajduje się w~fazie kontaktu lewej stopy (\texttt{LEFT\_STANCE}), prawej stopy (\texttt{RIGHT\_STANCE}) czy lotu (\texttt{FLIGHT}). Ręczna adnotacja ponad 12~tysięcy klatek byłaby czasochłonna i~obarczona subiektywną oceną momentu kontaktu, opracowałem więc algorytm automatycznego etykietowania na podstawie trajektorii keypointów stóp.

\subsection{Odrzucone podejście --- progowanie pozycji pięty}
\label{sec:labeling-rejected}

Pierwszym podejściem, które testowałem, było progowanie współrzędnej $y$ pięty: jeśli pięta znajduje się poniżej ustalonego progu (blisko podłoża), klatka klasyfikowana jest jako STANCE, w~przeciwnym razie jako FLIGHT. Podejście to okazało się niewystarczające z~trzech powodów:

\begin{enumerate}
    \item \textbf{Krótka faza lotu.} Przy 30~FPS faza lotu trwa zaledwie 2--3~klatki, co daje bardzo wąskie okno detekcji. Nawet niewielki błąd w~progu eliminuje te klatki --- przy progu 0{,}02 (w~znormalizowanej przestrzeni keypointów) algorytm klasyfikował ponad 70\% klatek jako FLIGHT.
    \item \textbf{Asymetria L/R z~kąta kamery.} Ujęcie boczne powoduje, że stopa bliższa kamery ma systematycznie inną wartość $y$ niż stopa dalsza. Stały próg nie radzi sobie z~tą asymetrią.
    \item \textbf{Utrata krótkich segmentów.} Filtr medianowy stosowany po progowaniu kasował jednoklatowe segmenty FLIGHT, zamiast je zachować.
\end{enumerate}

\subsection{Algorytm peak-based}
\label{sec:labeling-peak-based}

Finalne podejście opiera się na detekcji lokalnych maksimów (peaków) w~sygnale kontaktu stopy, co odpowiada momentom uderzenia stopy o~podłoże. Algorytm składa się z~pięciu kroków.

\subsubsection{Konstrukcja sygnału kontaktu}

Dla każdej stopy obliczałem sygnał kontaktu jako maksimum ze współrzędnych $y$ pięty i~czubka stopy:
\[
s_{\text{kontakt}}(t) = \max\bigl(y_{\text{heel}}(t),\; y_{\text{foot\_index}}(t)\bigr)
\]
Operacja $\max$ zapewnia, że sygnał obejmuje cały cykl kontaktu --- od momentu uderzenia piętą (heel strike, kiedy $y_{\text{heel}}$ jest maksymalne) do oderwania palców (toe-off, kiedy $y_{\text{foot\_index}}$ jest maksymalne). Stosowanie samego $y_{\text{heel}}$ pomijałoby końcową fazę kontaktu, a~samego $y_{\text{foot\_index}}$ --- początkową.
\newpage
\subsubsection{Detekcja foot strikes}

Momenty uderzenia stopy identyfikowałem osobno dla lewej i~prawej nogi za pomocą funkcji \texttt{scipy.signal.find\_peaks()} z~biblioteki SciPy~\cite{scipy}, z~dwoma parametrami:

\begin{itemize}
    \item \textbf{Minimalna odległość} (\texttt{distance}) --- wyznaczana dynamicznie ze wzoru:
    \[
    d_{\min} = \max\!\left(3,\;\left\lfloor \frac{\text{FPS} \cdot 60}{c_{\max}} \right\rfloor\right)
    \]
    gdzie $c_{\max} = 150$ to maksymalna spodziewana kadencja jednej stopy (kontaktów na minutę). Przy 30~FPS daje to $d_{\min} = 12$~klatek ($\approx$0{,}4~s), co zapobiega wykryciu dwóch peaków w~obrębie jednego kroku.
    \item \textbf{Minimalna prominencja} (\texttt{prominence} $= 0{,}03$) --- miara tego, jak bardzo peak wystaje ponad otaczający sygnał w~znormalizowanej przestrzeni keypointów $[0, 1]$. Próg ten odrzuca drobne oscylacje szumu, zachowując rzeczywiste kontakty stopy.
\end{itemize}

\subsubsection{Wymuszenie alternacji L--R}

W~prawidłowym biegu stopy lądują naprzemiennie: lewa, prawa, lewa, prawa. Algorytm scala peaki obu stóp w~jedną listę chronologiczną i~wymusza ten wzorzec: jeśli dwa kolejne peaki dotyczą tej samej stopy (np.~L--L), zachowuję ten z~wyższą prominencją, a~drugi odrzucam. Wynikiem jest ściśle naprzemienna sekwencja kontaktów L--R--L--R.

\subsubsection{Podział na fazy STANCE i~FLIGHT}

Między każdą parą kolejnych kontaktów algorytm wyznacza punkt przejścia --- klatkę, w~której wartość $\max(s_{\text{L}}, s_{\text{R}})$ osiąga minimum, co odpowiada momentowi, gdy obie stopy są najwyżej nad podłożem (szczyt fazy lotu). Każdy region między punktami przejścia dzielę na strefę centralną (STANCE) i~marginesy brzegowe (FLIGHT) za pomocą parametru $f_{\text{flight}} = 0{,}4$:
\[
m = \max\!\left(1,\;\left\lfloor \frac{l_{\text{region}} \cdot f_{\text{flight}}}{2} \right\rfloor\right)
\]
gdzie $l_{\text{region}}$ to długość regionu w~klatkach, a~$m$ to szerokość marginesu. Klatki w~pierwszych $m$ klatkach regionu oraz ostatnich $m$ klatkach oznaczam jako \texttt{FLIGHT}, a~klatki centralne jako odpowiednio \texttt{LEFT\_STANCE} lub \texttt{RIGHT\_STANCE}. \\Parametr $f_{\text{flight}} = 0{,}4$ oznacza, że około 40\% klatek między kolejnymi kontaktami przypisywanych jest do fazy lotu --- co jest zgodne z~typowymi proporcjami faz biegu rekreacyjnego \cite{novacheck1998}.

\subsubsection{Detekcja kierunku biegu i~filtr medianowy}

Algorytm automatycznie wykrywa kierunek biegu, porównując średnią pozycję $x$ nosa ze~średnią pozycją $x$ środka bioder w~całym nagraniu. Jeśli nos znajduje się na lewo od bioder (biegacz biegnie w~lewo w~kadrze), zamieniam etykiety \texttt{LEFT\_STANCE} $\leftrightarrow$ \texttt{RIGHT\_STANCE}, aby zachować anatomiczną poprawność niezależnie od orientacji nagrania.

Na zakończenie stosowałem filtr medianowy z rozmiarem okna $k = 3$ klatki --- dla każdej klatki filtr bierze ją samą oraz jedną klatkę przed i~jedną po, a~jako wynik przyjmuje etykietę występującą najczęściej w~tym trójkowym oknie, eliminując ewentualne pojedyncze klatki z~błędną etykietą. W~praktyce filtr ten zmieniał znikomą liczbę klatek --- algorytm peak-based generuje sekwencję wystarczająco czystą, że post-processing jest niemal zbędny. Wynikowe etykiety zapisywałem w~dwóch kolumnach pliku CSV z~keypointami: \texttt{phase\_auto} (surowa etykieta przed filtrem) i~\texttt{phase} (finalna etykieta po filtrze medianowym).

\section{Architektury klasyfikatorów}
\label{sec:architektury}

W~ramach pracy porównałem cztery architektury klasyfikatorów faz biegu, aby ocenić wpływ inżynierii cech oraz modelowania kontekstu czasowego na dokładność klasyfikacji. Modele różnią się zarówno algorytmem (Random Forest vs BiLSTM), jak i~reprezentacją danych wejściowych (surowe keypointy vs cechy inżynierowane).

\subsection{Random Forest --- baseline (RF\,v1)}
\label{sec:rf-v1}

Pierwszy model stanowi baseline --- punkt odniesienia weryfikujący, czy samo zadanie klasyfikacji faz biegu jest wykonalne na podstawie keypointów z~pojedynczej klatki. Random Forest~\cite{breiman2001rf} otrzymuje na wejściu wektor 132~surowych cech, czyli bezpośrednich, nieprzetworzonych wyjść z~MediaPipe Pose: $33 \text{ keypointy} \times 4 \text{ wartości } (x, y, z, visibility) = 132$. Surowe współrzędne zależą od pozycji biegacza w~kadrze i~proporcji jego ciała --- ta sama poza w~innym miejscu kadru daje inny wektor cech. Model musi samodzielnie nauczyć się ignorować te nieistotne różnice.

Hiperparametry (w~implementacji biblioteki scikit-learn~\cite{sklearn}): 300~drzew (\texttt{n\_estimators}), nieograniczona głębokość (\texttt{max\_depth=None}), zbalansowane wagi klas (parametr \texttt{class\_weight} ustawiony na \texttt{'balanced'}) kompensujące ewentualne nierówności w~rozkładzie faz.

\subsection{Random Forest z~cechami inżynierowanymi (RF\,v2)}
\label{sec:rf-v2}

Drugi model wykorzystuje 106~cech inżynierowanych, czyli wielkości obliczonych z~surowych keypointów w~celu uzyskania reprezentacji bardziej informatywnej i~niezależnej od warunków nagrania. Proces transformacji przebiega w~dwóch krokach:

\begin{enumerate}
    \item \textbf{Normalizacja do układu ciała.} Obliczałem punkt środkowy bioder (\textit{mid\_hip}) jako średnią pozycji lewego i~prawego biodra, a~następnie długość tułowia jako odległość euklidesową między \textit{mid\_hip} a~środkiem ramion (\textit{mid\_shoulder}). Każdy keypoint normalizowałem odejmując pozycję \textit{mid\_hip} i~dzieląc przez długość tułowia. Daje to $33 \times 3 = 99$~cech znormalizowanych, niezależnych od pozycji biegacza w~kadrze i~przeskalowanych względem proporcji ciała.
    \item \textbf{Kąty stawów.} Obliczałem 6~kątów stawów (po 3 na nogę): kąt kolana (biodro--kolano--kostka), kąt biodra (ramię--biodro--kolano) i~kąt kostki (kolano--kostka--czubek stopy), oraz pochylenie tułowia (kąt wektora \textit{mid\_hip}$\to$\textit{mid\_shoulder} względem pionu). Razem 7~dodatkowych cech, łącznie $99 + 7 = 106$.
\end{enumerate}

Hiperparametry Random Forest pozostawiłem identyczne jak w~RF\,v1 (300~drzew, brak ograniczenia głębokości, zbalansowane wagi klas), aby różnica w~wynikach odzwierciedlała wyłącznie wpływ inżynierii cech.

\subsection{BiLSTM --- model sekwencyjny}
\label{sec:bilstm}

Modele Random Forest klasyfikują każdą klatkę niezależnie, ignorując kontekst czasowy --- nie wiedzą, co działo się w~klatkach sąsiednich. Tymczasem fazy biegu tworzą uporządkowaną sekwencję (kontakt lewej stopy $\to$ lot $\to$ kontakt prawej stopy $\to$ lot), co sugeruje, że model uwzględniający kolejność klatek powinien osiągnąć lepsze wyniki.

\subsubsection{Architektura sieci}

Zastosowałem sieć typu LSTM (Long Short-Term Memory) \cite{hochreiter1997lstm} --- rodzaj sieci neuronowej rekurencyjnej zaprojektowanej do przetwarzania sekwencji. W~odróżnieniu od Random Forest, LSTM przetwarza ciąg kolejnych klatek i~na każdym kroku decyduje, które informacje z~wcześniejszych klatek zapamiętać, a~które zapomnieć. Dzięki temu model może uchwycić wzorce czasowe --- np.~że po fazie lotu zawsze następuje kontakt stopy.

Sieć jest dwukierunkowa (BiLSTM), co oznacza, że przetwarza sekwencję zarówno od~lewej do~prawej, jak i~od~prawej do~lewej. Wyniki obu kierunków są łączone, dzięki czemu dla każdej klatki model widzi kontekst zarówno z~przeszłości, jak i~z~przyszłości.

Na wejściu sieć otrzymuje okno 15~kolejnych klatek (po 106~cech inżynierowanych na klatkę) i~klasyfikuje fazę klatki środkowej -- 7. Okno 15~klatek przy 30~FPS obejmuje $\approx$ 0{,}5~sekundy, co odpowiada w~przybliżeniu jednemu pełnemu cyklowi biegowemu --- wystarczająco dużo kontekstu, by model widział przejścia między fazami. 

Cechy przed podaniem do sieci normalizowałem za pomocą \texttt{StandardScaler}: od~każdej cechy odejmowałem jej średnią i~dzieliłem przez odchylenie standardowe, tak aby wszystkie cechy miały porównywalny zakres wartości (średnia~$\approx 0$, odchylenie~$\approx 1$). Średnią i~odchylenie standardowe obliczałem wyłącznie na danych treningowych, a~następnie tymi samymi wartościami transformowałem dane walidacyjne i~testowe. Gdybym obliczył statystyki na całym zbiorze łącznie, model pośrednio ``poznałby'' rozkład danych testowych jeszcze przed testowaniem --- wyniki byłyby sztucznie zawyżone i~nie odzwierciedlałyby rzeczywistej zdolności modelu do generalizacji na nowe, nieznane dane. 

\subsubsection{Trening sieci}

Trening sieci neuronowej~\cite{goodfellow2016deep} polega na iteracyjnym dostrajaniu jej parametrów (wag), tak aby predykcje modelu były jak najbliższe prawdziwym etykietom. Sieć zaimplementowałem w~bibliotece PyTorch~\cite{pytorch}. Proces wymaga zdefiniowania trzech elementów:

\begin{itemize}
    \item \textbf{Funkcja straty} --- mierzy, jak bardzo predykcja modelu odbiega od prawdziwej etykiety. Zastosowałem \texttt{CrossEntropyLoss}, standardową funkcję straty dla klasyfikacji wieloklasowej, która karze model proporcjonalnie do pewności, z~jaką wskazuje błędną klasę. Ponieważ trzy fazy biegu występują w~zbiorze z~różną częstością, zastosowałem zbalansowane wagi klas --- fazy rzadsze mają proporcjonalnie wyższą wagę w~funkcji straty, co zapobiega faworyzowaniu klasy dominującej.
    \item \textbf{Optymalizator} --- algorytm, który na podstawie wartości straty koryguje wagi sieci. Wybrałem Adam \cite{kingma2015adam}, który automatycznie dostosowuje tempo uczenia osobno dla każdego parametru. Jest to najpowszechniej stosowany optymalizator w~głębokim uczeniu, ponieważ dobrze radzi sobie z~różnorodnymi zadaniami bez konieczności szczegółowego dostrajania.
    \item \textbf{Tempo uczenia} --- określa, jak duże korekty wag wykonuje optymalizator w~każdym kroku. Zbyt wysokie tempo powoduje niestabilność (model ``przeskakuje'' optimum), zbyt niskie --- bardzo wolną zbieżność.
\end{itemize}

\subsubsection{Dwie konfiguracje --- LSTM\,r1 i~LSTM\,r2}

Wytrenowałem dwie konfiguracje BiLSTM, aby zbadać, czy mniejszy model z~silniejszą regularyzacją lepiej generalizuje na nowe filmy. Regularyzacja to zbiorcze określenie na techniki ograniczające \textit{overfitting}, czyli przeuczenie --- sytuację, w~której model zbyt dokładnie zapamiętuje dane treningowe (włącznie z~ich szumem i~przypadkowymi wzorcami) i~osiąga wysoką dokładność na danych treningowych, ale znacząco gorszą na nowych, niewidzianych wcześniej danych.

\begin{itemize}
    \item \textbf{LSTM\,r1} --- większy model o~ukrytym wymiarze $h = 128$ na kierunek. Ukryty wymiar to rozmiar wektora pamięci wewnętrznej sieci, czyli ile liczb sieć wykorzystuje do zakodowania informacji o~dotychczasowej sekwencji na każdym kroku. Większa wartość pozwala uchwycić bardziej złożone wzorce, ale zwiększa ryzyko przeuczenia. Ponieważ sieć jest dwukierunkowa, każda klatka otrzymuje dwa wektory pamięci --- jeden z~przejścia ``od lewej do prawej'' i~drugi z~przejścia ``od prawej do lewej'' --- które są łączone (konkatenowane) w~jeden wektor o~wymiarze $2 \times 128 = 256$. Model ma łącznie 637\,699~parametrów, czyli tylu liczbowych wag, które sieć dostosowuje w~trakcie treningu. Tempo uczenia: $0{,}001$.
    \item \textbf{LSTM\,r2} --- mniejszy model: $h = 64$ (128~po połączeniu obu kierunków), czyli 187\,779~parametrów --- redukcja o~71\% względem LSTM\,r1. Niższe tempo uczenia ($0{,}0003$) sprawia, że model uczy się ostrożniej --- korekty wag w~każdym kroku są mniejsze.
\end{itemize}

Oba modele stosują dwa mechanizmy ograniczające przeuczenie:

\begin{itemize}
\item \textbf{Dropout} --- podczas treningu losowo wyłącza część neuronów. W~modelu LSTM\,r1 wyłączałem 30\% neuronów, a~w~LSTM\,r2 --- 40\%. Dzięki temu sieć nie może zbyt mocno polegać na pojedynczych neuronach i~jest mniej podatna na zapamiętywanie przypadkowych wzorców z~danych treningowych. Większy dropout w~LSTM\,r2 zastosowałem, aby silniej przeciwdziałać przeuczeniu.

\item \textbf{Weight decay} --- ogranicza nadmierny wzrost wag modelu. Można to rozumieć jako dodatkową karę za zbyt skomplikowane rozwiązania. W modelu LSTM\,r1 zastosowałem wartość \(10^{-5}\), a w LSTM\,r2 większą wartość \(10^{-4}\), aby mocniej ograniczyć przeuczenie. Dzięki temu LSTM\,r2 miał być prostszym i ostrożniej uczącym się modelem.

\end{itemize}


Trening kończył się automatycznie mechanizmem \textit{early stopping}: jeśli strata na zbiorze walidacyjnym nie poprawiała się przez 15~kolejnych epok, przerywałem trening i~zachowywałem wagi z~najlepszej epoki. Zapobiegało to sytuacji, w~której dłuższy trening pogarsza generalizację. LSTM\,r1 osiągnął najlepszą epokę jako 5.~z~20, LSTM\,r2 --- jako 3.~z~18, co wskazuje na szybką zbieżność obu modeli.

\subsection{Korekta proporcji obrazu (aspect ratio fix)}
\label{sec:aspect-fix}

W~trakcie eksperymentów zaobserwowałem, że modele z~cechami inżynierowanymi (RF\,v2, oba LSTM) osiągały systematycznie gorsze wyniki na filmie~10 (wideo pionowe, $608 \times 1080$~px). Przyczyną okazał się sposób, w~jaki MediaPipe normalizuje współrzędne keypointów.

MediaPipe zwraca współrzędne $x$ i~$y$ w~zakresie $[0, 1]$, niezależnie od rzeczywistych proporcji kadru. Oznacza to, że przesunięcie keypointu o~$\Delta x = 0{,}1$ w~filmie o~szerokości 608~pikseli odpowiada 60{,}8~pikseli, natomiast przesunięcie o~$\Delta y = 0{,}1$ przy wysokości 1080~pikseli odpowiada 108~pikseli --- prawie dwukrotnie więcej. Innymi słowy: ta sama zmiana w~znormalizowanych współrzędnych oznacza fizycznie inny dystans w~pionie niż w~poziomie. Problem ten ujawnia się przy obliczaniu odległości euklidesowych, np.~długości tułowia --- formuła traktuje $\Delta x$ i~$\Delta y$ jako porównywalne, podczas gdy w~rzeczywistości odpowiadają one różnej liczbie pikseli. Przy filmach o~standardowych proporcjach 16:9 (np.~$1920 \times 1080$) różnica jest umiarkowana, natomiast przy wideo pionowym ($608 \times 1080$, proporcje $\approx$~9:16) zniekształcenie jest znaczące.

Korektę zaimplementowałem jako krok preprocessingu przed obliczaniem cech inżynierowanych: mnożę współrzędne $x$ przez szerokość kadru, a~$y$ przez wysokość, przywracając fizycznie spójne jednostki (piksele). Dzięki temu $\Delta x$ i~$\Delta y$ wyrażają rzeczywiste odległości w~pikselach, a~obliczane na ich podstawie kąty i~odległości odpowiadają geometrii ciała biegacza, nie proporcjom kadru. Dla filmów o~standardowych proporcjach 16:9 korekta ma niewielki wpływ, natomiast dla wideo pionowego (film~10) poprawiła test accuracy modelu LSTM\,r1 na tym filmie z~75{,}8\% do~85{,}9\% (+10{,}1~pp). Łączna poprawa na całym zbiorze testowym wyniosła +3{,}8~pp (z~67{,}1\% do~70{,}9\%).

Tabela~\ref{tab:classifiers} podsumowuje architektury wszystkich czterech modeli.

\begin{table}[htbp]
\centering
\caption{Porównanie architektur klasyfikatorów faz biegu}
\label{tab:classifiers}
\small
\begin{tabular}{lcccc}
\hline
 & \textbf{RF\,v1} & \textbf{RF\,v2} & \textbf{LSTM\,r1} & \textbf{LSTM\,r2} \\
\hline
Algorytm       & Random Forest & Random Forest & BiLSTM & BiLSTM \\
Cechy           & 132 (surowe) & \multicolumn{3}{c}{106 (inżynierowane)} \\
Kontekst czasowy & 1~klatka & 1~klatka & 15~klatek & 15~klatek \\
Aspect fix      & nie & tak & tak & tak \\
Ukryty wymiar   & --- & --- & 128 & 64 \\
Dropout         & --- & --- & 0{,}3 & 0{,}4 \\
Parametry       & --- & --- & 637\,699 & 187\,779 \\
\hline
\end{tabular}
\end{table}

\section{Obliczanie współczynników biomechanicznych}
\label{sec:obliczanie-wspolczynnikow}

Po sklasyfikowaniu faz biegu pipeline oblicza 12~współczynników biomechanicznych opisujących technikę biegacza. Współczynniki podzieliłem na dwie grupy:

\begin{itemize}
    \item \textbf{Temporalne} (czasowe) --- oparte na czasie trwania poszczególnych faz biegu. Nazywam je temporalnymi, ponieważ ich wartość wynika z~tego, jak długo trwają kolejne fazy (kontakt, lot), a~nie z~tego, jaką pozycję ciało przyjmuje. Do ich obliczenia wystarczy sekwencja etykiet faz i~częstotliwość klatek wideo. 
    \item \textbf{Przestrzenne} --- oparte na geometrii keypointów w~konkretnych klatkach, czyli na tym, gdzie w~przestrzeni obrazu znajdują się poszczególne punkty ciała. 
\end{itemize}

W~typowym nagraniu biegacz wykonuje kilkanaście do kilkudziesięciu cykli biegowych. Każdy cykl daje jedną wartość danego współczynnika --- np.~jeden czas kontaktu lewej stopy, jeden czas lotu, jeden kąt kolana w~momencie kontaktu. Wynikowy współczynnik raportowałem jako średnią arytmetyczną $\pm$~odchylenie standardowe ze~wszystkich cykli w~nagraniu, co pozwala ocenić zarówno typową wartość, jak i~jej zmienność od cyklu do cyklu. Współczynniki zależne od strony ciała --- takie jak czas kontaktu czy kąt kolana --- obliczałem osobno dla lewej i~prawej nogi, uzyskując dwie niezależne wartości (np.~GCT$_L$ i~GCT$_R$), które następnie porównywałem wskaźnikiem symetrii.

\subsection{Segmentacja sekwencji faz}
\label{sec:segmentacja}

Klasyfikator faz przypisuje każdej klatce wideo jedną z~trzech etykiet: \texttt{LEFT\_STANCE}, \texttt{RIGHT\_STANCE} lub \texttt{FLIGHT}. Aby obliczyć współczynniki temporalne, muszę wiedzieć, jak długo trwa każda kolejna faza --- np.~ile klatek z~rzędu biegacz spędza w~fazie kontaktu lewej stopy, zanim przejdzie do fazy lotu. W~tym celu dzielę sekwencję etykiet na ciągłe fragmenty (segmenty) o~jednakowej etykiecie.

Algorytm przebiega sekwencję etykiet od pierwszej do ostatniej klatki: dopóki kolejne klatki mają tę samą etykietę, należą do tego samego segmentu. Gdy etykieta się zmienia, bieżący segment się kończy i~zaczyna nowy. Metodę tę określa się jako \textit{run-length encoding}. Przykładowo, sekwencja 12~klatek:

\begin{center}
\small
\texttt{L, L, L, F, F, R, R, R, R, F, F, L}
\end{center}

\noindent
zostaje podzielona na pięć segmentów: \texttt{LEFT\_STANCE} (3~klatki), \texttt{FLIGHT} (2~klatki), \texttt{RIGHT\_STANCE} (4~klatki), \texttt{FLIGHT} (2~klatki), \texttt{LEFT\_STANCE} (1~klatka). Dla każdego segmentu zapisuję fazę, indeks klatki początkowej i~końcowej oraz czas trwania obliczony ze~wzoru:
\[
t_{\text{segment}} = \frac{n_{\text{klatek}}}{\text{FPS}} \quad [\text{s}]
\]
Na~przykład segment \texttt{RIGHT\_STANCE} trwający 4~klatki przy 30~FPS ma czas trwania $4 / 30 \approx 0{,}133$~s, czyli 133~ms. Te czasy trwania poszczególnych segmentów stanowią podstawę do obliczenia współczynników temporalnych opisanych w~kolejnych sekcjach.

\subsection{Współczynniki temporalne}
\label{sec:temporal}

\subsubsection{Kadencja}

Kadencja określa liczbę kontaktów stóp z~podłożem na minutę i~jest jednym z~podstawowych wskaźników ekonomii biegu. Obliczałem ją jako:
\[
\text{kadencja} = \frac{n_{\text{kroków}}}{t_{\text{nagrania}}} \cdot 60 \quad [\text{spm}]
\]
gdzie $n_{\text{kroków}}$ to liczba przejść z~fazy FLIGHT do STANCE (lewej lub prawej), a~$t_{\text{nagrania}}$ to czas trwania analizowanego fragmentu w~sekundach. Krok definiowałem jako wejście w~fazę kontaktu, czyli przejście z~klatki niebędącej STANCE do klatki STANCE. Kadencja typowego biegacza rekreacyjnego mieści się w~zakresie 150--170~spm, a~biegaczy zaawansowanych w~zakresie 170--185~spm \cite{heiderscheit2011}.

\subsubsection{Czas kontaktu z~podłożem (GCT)}

Czas kontaktu z~podłożem (\textit{Ground Contact Time}, GCT) to czas, przez jaki stopa pozostaje w~kontakcie z~podłożem w~trakcie jednego kroku. Segment STANCE to ciągły fragment klatek, w~których stopa pozostaje na ziemi --- np.~6~kolejnych klatek z~etykietą \texttt{LEFT\_STANCE} stanowi jeden segment, czyli jeden pojedynczy kontakt lewej stopy z~podłożem. W~typowym 20-sekundowym nagraniu mieści się kilkanaście takich segmentów na nogę.

Obliczałem GCT osobno dla lewej i~prawej nogi jako średni czas trwania segmentów o~etykiecie odpowiednio \texttt{LEFT\_STANCE} i~\texttt{RIGHT\_STANCE}:
\[
\text{GCT} = \frac{1}{N} \sum_{i=1}^{N} t_{\text{stance},\,i} \quad [\text{ms}]
\]
gdzie $N$ to łączna liczba segmentów STANCE danej nogi w~nagraniu, indeks $i$ numeruje kolejne segmenty (od pierwszego do $N$-tego), a~$t_{\text{stance},\,i}$ to czas trwania $i$-tego segmentu w~milisekundach, obliczony z~liczby klatek segmentu i~FPS nagrania. Wzór ten oznacza po prostu średnią arytmetyczną wszystkich pojedynczych czasów kontaktu zarejestrowanych dla danej nogi. Typowy GCT dla biegu rekreacyjnego wynosi 200--280~ms --- krótszy GCT jest kojarzony z~lepszą ekonomią biegu \cite{souza2016}.

\subsubsection{Czas lotu}

Czas lotu  to czas, w~którym obie stopy znajdują się w~powietrzu. Obliczałem go analogicznie do GCT --- jako średni czas trwania segmentów FLIGHT. Typowy zakres to 80--150~ms; brak fazy lotu oznacza chód, nie bieg \cite{novacheck1998}.

\subsubsection{Czas cyklu i~długość kroku}

Czas cyklu to czas od kontaktu jednej stopy do następnego kontaktu \textit{tej samej} stopy, obejmujący pełen cykl ruchu: kontakt~$\to$~lot~$\to$~kontakt drugiej stopy~$\to$~lot~$\to$~powrót do kontaktu pierwszej stopy. Obliczałem go osobno dla lewej i~prawej nogi jako różnicę znaczników czasowych kolejnych wejść w~STANCE tej samej nogi.

Jeśli użytkownik podał prędkość bieżni $v$~[m/s], obliczałem długość kroku:
\[
\text{stride} = v \cdot t_{\text{cykl}} \quad [\text{m}]
\]
Formuła ta wykorzystuje fakt, że na bieżni mechanicznej biegacz porusza się ze~stałą prędkością wymuszoną przez taśmę bieżni. Prędkość nie jest możliwa do wyznaczenia z~samego nagrania wideo --- wymaga podania przez użytkownika. 

\subsubsection{Duty factor}

Duty factor to stosunek czasu kontaktu stopy z~podłożem do czasu cyklu.
\[
\text{DF} = \frac{\text{GCT}}{t_{\text{cykl}}}
\]
Wartość poniżej~0{,}5 oznacza, że stopa spędza mniej niż połowę cyklu na podłożu, co jest charakterystyczne dla biegu. Wartość powyżej~0{,}5 oznacza, że stopa jest dłużej na ziemi niż w~powietrzu --- jest to cecha chodu, nie biegu. Typowy zakres dla biegu rekreacyjnego to 0{,}35--0{,}45, dla sprintów 0{,}22--0{,}30 \cite{souza2016}.

\subsection{Współczynniki przestrzenne}
\label{sec:spatial}

Współczynniki przestrzenne obliczałem na podstawie geometrii keypointów --- czyli ich położenia w~przestrzeni obrazu. Wykorzystywałem keypointy po filtracji filtrem Savitzky-Golay, nczyli \textit{wygładzone keypointy} --- są to surowe pozycje z~MediaPipe pozbawione szybkich, losowych fluktuacji (jittera), ale zachowujące rzeczywisty kształt ruchu ciała.

Część współczynników obliczałem w~pojedynczych klatkach --- np.~kąt kolana wyznaczany w~konkretnej klatce wejścia stopy w~fazę STANCE, dający jedną wartość na każdy kontakt. Część obliczałem z~całych faz --- np.~średni kąt kolana w~ciągu wszystkich klatek fazy STANCE, dający jedną uśrednioną wartość na każdy segment. Wartości raportowałem jako średnią~$\pm$~odchylenie standardowe ze~wszystkich cykli biegowych, z~rozbiciem na fazy (\texttt{LEFT\_STANCE}, \texttt{RIGHT\_STANCE}, \texttt{FLIGHT}) tam, gdzie miało to sens biomechaniczny.

\subsubsection{Kąty stawów}

Obliczałem sześć kątów stawów (po trzy na nogę) za pomocą funkcji trygonometrycznej:
\[
\alpha = \arccos\!\left(\frac{\vec{u} \cdot \vec{v}}{\|\vec{u}\| \cdot \|\vec{v}\|}\right) \cdot \frac{180^\circ}{\pi}
\]
gdzie $\vec{u}$ i~$\vec{v}$ to wektory wychodzące od keypointu stawu do dwóch sąsiednich keypointów:

\begin{itemize}
    \item \textbf{Kąt kolana} --- wektory: biodro~$\to$~kolano i~kolano~$\to$~kostka. Noga wyprostowana $\approx$~$170$--$180^\circ$, ugięta $\approx$~$90$--$120^\circ$.
    \item \textbf{Kąt biodra} --- wektory: ramię~$\to$~biodro i~biodro~$\to$~kolano. Opisuje rozwarcie uda względem tułowia.
    \item \textbf{Kąt kostki} --- wektory: kolano~$\to$~kostka i~kostka~$\to$~czubek stopy. Opisuje zgięcie stopy.
\end{itemize}

W~biomechanice biegu wyróżnia się dwa stany danej nogi w~obrębie cyklu:

\begin{itemize}
    \item \textbf{Noga oporowa} --- noga, której stopa aktualnie dotyka podłoża. Jest to noga utrzymująca ciężar ciała i~przejmująca siły reakcji podłoża. Odpowiada fazie \texttt{LEFT\_STANCE} lub \texttt{RIGHT\_STANCE} sklasyfikowanej przez model.
    \item \textbf{Noga w~fazie wahadła} --- noga uniesiona, przesuwająca się do przodu, aby przygotować się do kolejnego kontaktu. W~tym czasie druga noga jest albo w~kontakcie (STANCE) albo obie nogi są w~powietrzu (FLIGHT).
\end{itemize}

W~dowolnym momencie biegu jedna noga jest oporowa, a~druga w~wahadle --- nigdy nie obie naraz w~tym samym stanie. Kąty kolan obu nóg zmieniają się w~zależności od tego, w~którym stanie się znajdują --- np.~kąt kolana nogi oporowej (w~fazie STANCE) wynosi typowo $155$--$175^\circ$ (noga prawie wyprostowana, by amortyzować uderzenie), a~tej samej nogi w~fazie wahadła spada do~$100$--$140^\circ$ (silne zgięcie, by skrócić długość wahadła i~przyspieszyć przeniesienie nogi do przodu) \cite{novacheck1998}.

\subsubsection{Kąt kolana w~momencie kontaktu}

Oprócz kątów uśrednionych w~ramach całej fazy, obliczałem kąt kolana w~konkretnej klatce wejścia stopy w~fazę STANCE. Wartość ta jest wskaźnikiem overstridingu: kąt $160$--$175^\circ$ oznacza prawidłowe, lekko ugięte kolano absorbujące siłę uderzenia, natomiast kąt powyżej~$175^\circ$ wskazuje, że stopa ląduje daleko przed środkiem ciężkości, co zwiększa siły hamowania i~ryzyko kontuzji kolana \cite{heiderscheit2011}.

\subsubsection{Pochylenie tułowia}

Pochylenie tułowia obliczałem jako kąt między wektorem tułowia (od~środka bioder do~środka ramion) a~kierunkiem pionu w~płaszczyźnie obrazu. Wartości powyżej~$0^\circ$ oznaczają pochylenie do przodu. Zalecany zakres to $5$--$15^\circ$ --- lekkie pochylenie z~kostek wykorzystuje siłę ciężkości do napędu, natomiast pochylenie powyżej~$20^\circ$ nadmiernie obciąża dolny odcinek kręgosłupa \cite{novacheck1998}.

\subsubsection{Oscylacja pionowa}

Oscylacja pionowa  mierzy amplitudę ruchu bioder w~pionie w~obrębie jednego cyklu biegowego:
\[
\Delta y = \max(y_{\text{mid\_hip}}) - \min(y_{\text{mid\_hip}})
\]
obliczaną w~oknie od jednego wejścia w~\texttt{LEFT\_STANCE} do~następnego. Wartość surową normalizowałem względem średniej długości tułowia w~danym cyklu, dzięki czemu uzyskałem bezwymiarowy wskaźnik niezależny od odległości kamery i~rozdzielczości nagrania.

Długość tułowia w~pojedynczej klatce wyznaczałem jako odległość euklidesową między środkiem bioder a~środkiem ramion:
\[
l_{\text{torso}}(t) = \left\| \mathbf{p}_{\text{mid\_shoulder}}(t) - \mathbf{p}_{\text{mid\_hip}}(t) \right\|
\]
Następnie dla całego cyklu biegowego obliczałem średnią długość tułowia:
\[
L_{\text{torso}} = \frac{1}{N} \sum_{t=1}^{N} l_{\text{torso}}(t)
\]
gdzie $N$ oznacza liczbę klatek w~analizowanym cyklu biegowym.\\ Znormalizowaną oscylację pionową obliczałem jako:
\[
\text{VO}_{\text{norm}} = \frac{\Delta y}{L_{\text{torso}}}
\]

Uśrednianie długości tułowia w~obrębie cyklu jest konieczne, ponieważ wartość ta drobno waha się z~klatki na klatkę, zarówno z~powodu rzeczywistego ruchu ciała, jak i~resztkowego szumu MediaPipe. Wzięcie średniej daje stabilną wartość odniesienia.

Typowy zakres wynosi 0{,}12--0{,}16 długości tułowia ($\approx$~6--8~cm) dla dobrego biegacza i~do~0{,}24 ($\approx$~12~cm) dla biegacza rekreacyjnego. Wyższe wartości oznaczają nadmierne ``podskakiwanie'', czyli marnowanie energii na ruch pionowy zamiast na przemieszczanie się do przodu \cite{novacheck1998}.

\subsubsection{Wzorzec lądowania stopy}

Wzorzec lądowania stopy opisuje, która część stopy pierwsza dotyka podłoża w~momencie kontaktu. Klasyfikowałem go na podstawie kąta wektora od~pięty (\texttt{HEEL}) do~czubka stopy (\texttt{FOOT\_INDEX}) w~klatce wejścia w~fazę STANCE:
\[
\theta = \arctan2(-\Delta y,\; \Delta x)
\]
gdzie:
\begin{itemize}
    \item $\Delta x = x_{\text{foot\_index}} - x_{\text{heel}}$ --- różnica współrzędnych poziomych czubka stopy i~pięty (na ile pikseli czubek jest dalej od pięty w~poziomie),
    \item $\Delta y = y_{\text{foot\_index}} - y_{\text{heel}}$ --- różnica współrzędnych pionowych (na ile pikseli czubek jest wyżej lub niżej niż pięta). Znak minus przed $\Delta y$ we wzorze odwraca oś $y$, ponieważ w~MediaPipe oś pionowa rośnie w~dół (0~na górze, 1~na dole kadru), a~chcę mieć konwencję matematyczną (oś $y$ rosnąca w~górę),
    \item $\theta$ --- wynikowy kąt stopy w~stopniach. Geometrycznie reprezentuje on nachylenie stopy względem poziomu w~momencie kontaktu z~podłożem: $\theta > 0$ oznacza, że czubek stopy jest wyżej niż pięta (pięta nisko, pierwsza dotyka ziemi); $\theta < 0$ oznacza, że pięta jest wyżej niż czubek (przód stopy nisko, pierwszy dotyka ziemi); $\theta \approx 0$ to stopa równoległa do podłoża.
\end{itemize}

Na podstawie wartości $\theta$ klasyfikowałem wzorzec lądowania: $\theta > 5^\circ$ oznacza lądowanie na pięcie, $\theta < -5^\circ$ --- lądowanie na przodzie stopy, wartości pomiędzy --- lądowanie na śródstopiu. Progi $\pm  5^\circ$ przyjąłem jako wąską strefę neutralną wokół poziomu. Dominujący wzorzec dla nogi wyznaczałem na podstawie najczęściej występującej kategorii ze~wszystkich kontaktów tej nogi w~nagraniu \cite{daoud2012}.

Współczynnik ten jest wrażliwy na perspektywę kamery --- przy ujęciu innym niż ściśle boczne (np.~wideo pionowe) kąty stopy są zniekształcone. Dlatego zaimplementowałem automatyczną flagę \texttt{low\_confidence}: jeśli średni kąt stopy przekraczał $45^\circ$ w~wartości bezwzględnej, wzorzec lądowania oznaczałem jako niewiarygodny. 

\subsection{Symetria lewo-prawo}
\label{sec:symmetry}

W~prawidłowo wykonywanym biegu lewa i~prawa noga powinny pracować podobnie --- mieć zbliżone czasy kontaktu, zbliżone kąty stawów, zbliżoną długość cyklu. Różnice mogą sygnalizować niestabilność biegu, dysbalans mięśniowy lub inne anomalię biegu. Aby ilościowo opisać te różnice, dla każdego współczynnika obliczanego osobno dla lewej i~prawej strony wyznaczałem wskaźnik symetrii Robinsona (\textit{Symmetry Index}, SI), zaproponowany pierwotnie do oceny chodu po manipulacjach chiropraktycznych \cite{robinson1987}, a~obecnie szeroko stosowany w~analizie biegu:
\[
\text{SI} = \frac{200 \cdot |L - R|}{L + R} \quad [\%]
\]
gdzie $L$ i~$R$ to wartości danego współczynnika dla lewej i~prawej nogi. Sposób działania wzoru:

\begin{itemize}
    \item $|L - R|$ to bezwzględna różnica między stroną lewą a~prawą --- mierzy, „o~ile" się różnią,
    \item $L + R$ to suma obu wartości --- służy do normalizacji, dzięki czemu wynik nie zależy od skali współczynnika (SI~=~5\% oznacza tę samą względną asymetrię niezależnie od tego, czy mierzymy GCT w~milisekundach, czy kąt kolana w~stopniach),
    \item mnożnik 200 zapewnia, że wynik wyrażony jest w~procentach względem średniej z~lewej i~prawej (czyli $(L+R)/2$).
\end{itemize}

Dla przykładu: GCT$_L = 250$~ms, GCT$_R = 240$~ms $\Rightarrow$ SI $= 200 \cdot 10 / 490 \approx 4{,}1\%$. SI~=~0\% oznacza idealną symetrię. Na podstawie literatury przyjąłem trzy progi interpretacyjne: poniżej~5\% --- norma, 5--10\% --- łagodna asymetria wymagająca monitorowania, powyżej~10\% --- asymetria znacząca, potencjalnie wskazująca na niestabilność biegu lub kompensacje ruchu po kontuzji.

\subsubsection{Dla jakich współczynników obliczam SI}

Wskaźnik symetrii obliczałem dla pięciu par współczynników:
\begin{itemize}
    \item GCT (czas kontaktu) --- L vs~R,
    \item Czas cyklu --- L vs~R,
    \item Duty factor --- L vs~R,
    \item Kąt kolana w~fazie STANCE --- L vs~R,
    \item Kąt kostki w~fazie STANCE --- L vs~R.
\end{itemize}

Dodatkowo porównywałem wzorzec lądowania stopy (heel / midfoot / forefoot) --- sprawdzając, czy obie stopy lądują tym samym wzorcem. Tu nie stosowałem wzoru SI, ponieważ wzorzec jest zmienną kategoryczną (heel/midfoot/forefoot), a~nie liczbową --- zamiast tego zwracałem prostą flagę „spójność: tak/nie".

\subsubsection{Dlaczego dla każdego współczynnika osobno, a~nie jedna miara „symetrii nogi"}

Naturalne pytanie brzmi: czy nie wystarczyłoby obliczyć \textit{jednej} ogólnej miary symetrii między lewą a~prawą nogą? Odpowiedź brzmi: nie, ponieważ asymetria może manifestować się w~różny sposób i~różne typy asymetrii mają różne implikacje biomechaniczne. Biegacz może na przykład mieć:

\begin{itemize}
    \item \textbf{Symetryczny GCT, ale asymetryczne kąty kolan}
    \item \textbf{Asymetryczny GCT, ale symetryczne kąty} 
    \item \textbf{Asymetryczny duty factor, symetryczny GCT i~czas cyklu}
\end{itemize}

Pojedyncza ogólna metryka zatraciłaby tę informację diagnostyczną. Dla każdego współczynnika SI mówi coś innego o~biomechanice biegacza, dlatego warto obliczać je osobno.

Pod względem kosztu obliczeniowego nie jest to znaczące obciążenie: dla każdej pary wartości $L$ i~$R$ (które i~tak są już wyliczone w~ramach obliczeń per-noga) zastosowanie wzoru SI to jedna operacja do obliczenia.
\section{Silnik reguł rekomendacji}
\label{sec:silnik-regul}

Ostatnim etapem pipeline'u jest generowanie rekomendacji treningowych na podstawie obliczonych współczynników. Rekomendacje opierają się na regułach zakodowanych ręcznie na podstawie literatury biomechanicznej --- nie są uczone z~danych, lecz odzwierciedlają ustalone w~literaturze zakresy referencyjne i~wzorce ryzyka. Każda reguła analizuje jeden aspekt techniki biegu (np.~kadencję, czas kontaktu, symetrię) i~zwraca rekomendację z~przypisanym poziomem ważności oraz źródłem literaturowym.

Architekturę silnika, pełny katalog reguł, reguły łączone oraz źródła literaturowe progów opisuję w~całości w~rozdziale~\ref{ch:rekomendacje}, poświęconym systemowi rekomendacji. W~tym miejscu sygnalizuję jedynie, że silnik reguł domyka pipeline --- przekształca surowe współczynniki w~konkretne, opatrzone źródłami wskazówki treningowe.

\section{Podział danych na zbiory treningowy, walidacyjny i~testowy}
\label{sec:splity}

Zbiór 15~nagrań podzieliłem na trzy rozłączne podzbiory: treningowy (10~filmów, $\approx$~9\,800~klatek), walidacyjny (2~filmy, $\approx$~740~klatek) i~testowy (3~filmy, $\approx$~1\,500~klatek). Podział wykonałem na poziomie całych filmów, czyli wszystkie klatki danego nagrania trafiają do tego samego podzbioru --- nie dzielę pojedynczego filmu między zbiory. Dzięki temu model nie widzi w~treningu żadnych klatek z~filmów testowych, co daje uczciwą ocenę zdolności do generalizacji na nowe nagrania.

\subsection{Przypisanie filmów do zbiorów}
\label{sec:splity-przypisanie}

Tabela~\ref{tab:splits} przedstawia przydział filmów do zbiorów wraz z~rozkładem klas.

\begin{table}[htbp]
\centering
\caption{Podział filmów na zbiory treningowy, walidacyjny i~testowy}
\label{tab:splits}
\small
\begin{tabular}{lccccc}
\hline
\textbf{Zbiór} & \textbf{Filmów} & \textbf{Klatek} & \textbf{L\_STANCE} & \textbf{R\_STANCE} & \textbf{FLIGHT} \\
\hline
Treningowy  & 10 & 9\,779 & 3\,334 (34{,}1\%) & 3\,105 (31{,}7\%) & 3\,340 (34{,}2\%) \\
Walidacyjny &  2 &    738 &    261 (35{,}4\%) &    253 (34{,}3\%) &    224 (30{,}4\%) \\
Testowy     &  3 & 1\,490 &    525 (35{,}2\%) &    489 (32{,}8\%) &    476 (31{,}9\%) \\
\hline
\end{tabular}
\end{table}

Trzy klasy faz biegu rozkładają się w~zbiorze treningowym w~przybliżeniu równomiernie (31{,}7--34{,}2\%), co minimalizuje ryzyko faworyzowania klasy dominującej. Proporcje w~zbiorach walidacyjnym i~testowym są zbliżone do treningowych.

Do zbioru testowego celowo wybrałem trzy nagrania o~zróżnicowanym charakterze:
\begin{itemize}
    \item film~2 --- bieg w~stałym tempie 13~km/h, niska rozdzielczość ($360 \times 450$~px),
    \item film~9 (segment~1) --- bieg zmiennym tempem (0{,}8--3{,}5~m/s),
    \item film~10 --- wyższe tempo.
\end{itemize}

\section{Metryki ewaluacji}
\label{sec:metryki}

Do oceny jakości klasyfikatorów faz biegu zastosowałem trzy metryki: dokładność ogólną (\textit{accuracy}), makro-uśrednioną miarę F1 (\textit{macro F1}) oraz macierz pomyłek (\textit{confusion matrix}). Wszystkie metryki obliczałem na zbiorze testowym (3~filmy, $\approx$~1\,500~klatek dla Random Forest, $\approx$~1\,450~klatek dla BiLSTM --- różnica wynika z~tego, że BiLSTM operuje na oknach 15~kolejnych klatek, więc kilkanaście klatek z~początku i~końca każdego filmu nie ma pełnego kontekstu i~jest pomijanych).

\subsection{Dokładność ogólna (accuracy)}
\label{sec:metryki-accuracy}

Dokładność ogólna to odsetek poprawnie sklasyfikowanych próbek:
\[
\text{accuracy} = \frac{\text{liczba poprawnych predykcji}}{\text{liczba wszystkich predykcji}}
\]
Metryka ta jest intuicyjna, ale przy niezbalansowanych klasach może być myląca --- model przypisujący wszystkim próbkom klasę dominującą uzyskałby wysoką dokładność bez rzeczywistej zdolności klasyfikacji. W~moim zbiorze klasy są przybliżeniu zbalansowane ($\approx$~33\% każda), więc accuracy stanowi rozsądną główną metrykę.

\subsection{Precyzja, czułość i~miara F1}
\label{sec:metryki-f1}

Dla każdej z~trzech klas obliczałem:

\begin{itemize}
    \item \textbf{Precyzja} (\textit{precision}) --- odsetek poprawnych predykcji wśród wszystkich próbek, którym model przypisał daną klasę. Odpowiada na pytanie: ``Gdy model mówi FLIGHT, jak często ma rację?''
    \[
    \text{precision}_c = \frac{TP_c}{TP_c + FP_c}
    \]
    \item \textbf{Czułość} (\textit{recall}) --- odsetek poprawnie wykrytych próbek wśród wszystkich próbek rzeczywiście należących do danej klasy. Odpowiada na pytanie: ``Jaki odsetek prawdziwych klatek FLIGHT model wykrył?''
    \[
    \text{recall}_c = \frac{TP_c}{TP_c + FN_c}
    \]
    \item \textbf{Miara F1} --- średnia harmoniczna precyzji i~czułości:
    \[
    F1_c = 2 \cdot \frac{\text{precision}_c \cdot \text{recall}_c}{\text{precision}_c + \text{recall}_c}
    \]
\end{itemize}

\noindent
gdzie $TP_c$ to liczba \textit{true positives} klasy $c$ (poprawnie rozpoznane próbki tej klasy), $FP_c$ --- \textit{false positives} (próbki innych klas błędnie przypisane do $c$), a~$FN_c$ --- \textit{false negatives} (próbki klasy $c$ błędnie przypisane do innej).

Miara F1 zdefiniowana powyżej jest \textit{per-klasa} --- dla problemu klasyfikacji trzech faz biegu mam więc trzy osobne wartości F1: jedną dla \texttt{FLIGHT}, jedną dla \texttt{LEFT\_STANCE} i~jedną dla \texttt{RIGHT\_STANCE}. Aby uzyskać jedną liczbę porównawczą między modelami, potrzebowałem sposobu na zagregowanie tych trzech wartości w~jedną.

Zastosowałem \textbf{makro-uśrednione F1} (\textit{macro F1}), czyli zwykłą średnią arytmetyczną miar F1 z~poszczególnych klas:
\[
F1_{\text{macro}} = \frac{1}{3} \sum_{c=1}^{3} F1_c
\]
gdzie indeks $c$ przebiega trzy klasy: \texttt{FLIGHT}, \texttt{LEFT\_STANCE}, \texttt{RIGHT\_STANCE}.

Słowo „makro" odróżnia ten sposób uśredniania od alternatywnego podejścia zwanego \textit{micro F1}, w~którym najpierw sumuje się surowe liczby TP, FP i~FN ze wszystkich klas, a~dopiero potem oblicza pojedynczą wartość F1. Różnica między tymi podejściami sprowadza się do tego, czy klasy są traktowane jednakowo:

\begin{itemize}
    \item \textbf{macro F1} --- każda klasa wnosi równy wkład do końcowego wyniku, niezależnie od tego, ile próbek danej klasy znajduje się w~zbiorze (rzadka klasa ma taki sam wpływ jak częsta),
    \item \textbf{micro F1} --- klasy są ważone proporcjonalnie do liczby próbek (klasa z~1000~próbek waży 10$\times$ więcej niż klasa z~100~próbek). Dla klasyfikacji wieloklasowej z~jedną prawdziwą etykietą na próbkę micro F1 jest równe dokładności ogólnej (\textit{accuracy}).
\end{itemize}

Makro-uśrednioną miarę F1 wybrałem, ponieważ fazy biegu mogą mieć różną liczność w~zbiorze testowym. Sama dokładność ogólna mogłaby zawyżać ocenę modelu, jeśli dobrze rozpoznawałby fazy częste, a~popełniał więcej błędów w~klasach rzadszych lub trudniejszych do rozróżnienia. Macro F1, traktując każdą fazę z~równą wagą, lepiej odzwierciedla, czy klasyfikator działa stabilnie dla wszystkich faz cyklu biegowego, a~nie tylko dla klas dominujących liczbowo.
\newpage
\subsection{Macierz pomyłek}
\label{sec:metryki-confusion}

Macierz pomyłek to tablica $3 \times 3$, w~której wiersz $i$ i~kolumna $j$ zawiera liczbę próbek klasy $i$ (prawdziwej), które model przypisał do klasy $j$ (predykowanej). Elementy na przekątnej odpowiadają poprawnym klasyfikacjom, a~elementy poza przekątną --- pomyłkom. Macierz pozwala zidentyfikować, które pary klas model najczęściej myli.

W~kontekście faz biegu szczególnie istotne są dwa typy pomyłek:
\begin{itemize}
    \item \textbf{Pomyłka L$\leftrightarrow$R} (zamiana lewej i~prawej nogi) --- wpływa na współczynniki obliczane osobno per noga (GCT L/R, symetria), ale nie na współczynniki agregowane (kadencja, średni GCT),
    \item \textbf{Pomyłka STANCE$\leftrightarrow$FLIGHT} (zamiana kontaktu z~lotem) --- wpływa na wszystkie współczynniki temporalne, ponieważ zmienia długość segmentów faz.
\end{itemize}


\chapter{Klasyfikator faz biegu --- wyniki i~wybór modelu}
\label{ch:klasyfikator}

W~tym rozdziale przedstawiam wyniki porównania czterech architektur klasyfikatorów faz biegu opisanych w~sekcji~\ref{sec:architektury}. Omawiam dokładność ogólną, strukturę błędów, wpływ inżynierii cech i~korekty proporcji obrazu, a~na końcu uzasadniam wybór modelu głównego używanego w~dalszych etapach pipeline'u.

\section{Porównanie czterech modeli}
\label{sec:porownanie-modeli}

Tabela~\ref{tab:model-results} zestawia wyniki wszystkich czterech modeli na zbiorze walidacyjnym i~testowym. Wszystkie modele korzystające z~cech inżynierowanych stosowały korektę proporcji obrazu. Wartości F1 w~tabeli to makro-uśrednione miary F1. Kolumna ``Różnica'' przedstawia różnicę dokładności między zbiorem walidacyjnym a~testowym (val acc $-$ test acc) i~służy jako miara generalizacji modelu.

Generalizacja to zdolność modelu do osiągania dobrych wyników nie tylko na danych, na których się uczył, ale także na nowych, wcześniej niewidzianych przykładach. W~moim przypadku ``nowe przykłady'' to nagrania z~filmów testowych, których model nigdy nie widział w~trakcie treningu --- inni biegacze, inne kąty kamery, inne rozdzielczości. Model, który dobrze generalizuje, radzi sobie z~takimi nowymi nagraniami niemal tak samo dobrze, jak z~tymi, na których się uczył. Model słabo generalizujący osiąga wysokie wyniki tylko na danych treningowych, a~na nowych --- znacznie gorsze, bo zapamiętał konkretne wzorce z~treningu zamiast nauczyć się ogólnych reguł rozróżniania faz biegu. Im większa różnica między dokładnością walidacyjną a~testową, tym gorzej model generalizuje.

\begin{table}[htbp]
\centering
\caption{Wyniki klasyfikatorów faz biegu na zbiorze walidacyjnym i~testowym}
\label{tab:model-results}
\small
\begin{tabular}{lcccccc}
\hline
\textbf{Model} & \textbf{Val acc} & \textbf{Val F1} & \textbf{Test acc} & \textbf{Test F1} & \textbf{Różnica} \\
\hline
RF\,v1 (132 surowe)         & 82{,}0\% & 0{,}816 & 62{,}7\% & 0{,}617 & 19{,}3~pp \\
RF\,v2 (106 inżynierowane)  & 80{,}5\% & 0{,}803 & 65{,}7\% & 0{,}650 & 14{,}8~pp \\
LSTM\,r1 ($h$=128)   & 87{,}2\% & 0{,}869 & \textbf{70{,}9\%} & \textbf{0{,}709} & 16{,}3~pp \\
LSTM\,r2 ($h$=64)    & 85{,}0\% & 0{,}845 & 68{,}2\% & 0{,}681 & 16{,}8~pp \\
\hline
\end{tabular}
\end{table}

Z~porównania wynikają trzy główne obserwacje:

\begin{enumerate}
    \item \textbf{Inżynieria cech poprawia generalizację.} RF\,v2, który korzysta z~cech inżynierowanych --- czyli z~keypointów znormalizowanych do układu ciała biegacza wraz z~obliczonymi kątami stawów --- osiągnął na teście o~3{,}0~pp wyższą dokładność niż RF\,v1 operujący na surowych keypointach (65{,}7\% vs 62{,}7\%), mimo niższej dokładności na walidacji (80{,}5\% vs 82{,}0\%). Mniejsza różnica między walidacją a~testem (14{,}8~pp vs 19{,}3~pp) wskazuje, że taka reprezentacja lepiej generalizuje na nowych biegaczy: kąt kolana czy znormalizowana pozycja stopy względem bioder przyjmują podobne wartości niezależnie od tego, gdzie w~kadrze znajduje się biegacz, jak duża jest jego sylwetka i~jakie ma proporcje ciała --- podczas gdy surowe współrzędne $x$, $y$ keypointów zmieniają się przy każdym z~tych czynników.

    \item \textbf{Kontekst czasowy daje istotną poprawę.} Oba modele LSTM, które przetwarzają okno 15~klatek zamiast pojedynczej klatki, osiągnęły wyższe wyniki testowe niż oba modele Random Forest. Najlepszy LSTM\,r1 (70{,}9\%) przewyższył najlepszy RF\,v2 (65{,}7\%) o~5{,}2~pp, co potwierdza założenie, że uwzględnianie kontekstu czasowego między klatkami pozwala wyraźnie poprawić jakość klasyfikacji faz biegu.

    \item \textbf{Większy model lepszy od mniejszego.} LSTM\,r1 (128~neuronów, 637\,699~parametrów) osiągnął o~2{,}7~pp wyższą test accuracy niż LSTM\,r2 (64~neurony, 187\,779~parametrów). Silniejsza regularyzacja w~LSTM\,r2 (dropout~0{,}4, weight decay~$10^{-4}$) nie zrekompensowała ograniczonej pojemności modelu.
\end{enumerate}

Wszystkie modele wykazują znaczną różnicę między dokładnością na walidacji a~na teście (14{,}8--19{,}3~pp), co wynika z~tego, że zbiór testowy zawiera biegaczy nieobecnych w~zbiorze treningowym, biegających w~innych warunkach (rozdzielczość, tempo, orientacja kamery). Jest to oczekiwane zachowanie --- modele muszą generalizować na nowe osoby i~warunki nagrania.
\newpage
\section{Macierz pomyłek i~analiza błędów}
\label{sec:confusion-matrices}

Tabela~\ref{tab:confusion-lstm} przedstawia macierz pomyłek modelu LSTM\,r1 na zbiorze testowym ($\approx$~1\,450~klatek). Wiersze odpowiadają prawdziwym klasom, kolumny --- klasom przewidywanym przez model.

\begin{table}[htbp]
\centering
\caption{Macierz pomyłek LSTM\,r1 na zbiorze testowym}
\label{tab:confusion-lstm}
\begin{tabular}{l|ccc|c}
 & \textbf{FLIGHT} & \textbf{L\_STANCE} & \textbf{R\_STANCE} & \textbf{Czułość} \\
\hline
\textbf{FLIGHT}      & 321 &  79 &  76 & 67{,}4\% \\
\textbf{L\_STANCE}   & 110 & 333 &  54 & 67{,}0\% \\
\textbf{R\_STANCE}   &  82 &  21 & 372 & 78{,}3\% \\
\hline
\textbf{Precyzja}    & 62{,}6\% & 76{,}9\% & 74{,}1\% & \\
\end{tabular}
\end{table}

Najwyższą czułość (recall) model osiąga dla klasy \texttt{RIGHT\_STANCE} (78{,}3\%), najniższą --- dla \texttt{LEFT\_STANCE} (67{,}0\%). Klasa \texttt{FLIGHT} jest najczęściej mylona z~obiema klasami STANCE: 79~klatek lotu model błędnie przypisał do~\texttt{LEFT\_STANCE}, a~76 do~\texttt{RIGHT\_STANCE}. Ta asymetria jest spodziewana --- faza lotu trwa zaledwie 2--3~klatki przy 30~FPS, co oznacza, że klatki graniczne (na przejściu STANCE$\to$FLIGHT) mogą wyglądać niejednoznacznie.

\subsection{Struktura błędów: L$\leftrightarrow$R vs STANCE$\leftrightarrow$FLIGHT}

Nie wszystkie błędy klasyfikacji mają jednakowe konsekwencje dla obliczanych współczynników. Pomyłki \textbf{L$\leftrightarrow$R} (zamiana lewej i~prawej nogi) bezpośrednio rujnują wskaźniki symetrii i~współczynniki obliczane osobno per noga (GCT~L/R, duty factor~L/R). Pomyłki \textbf{STANCE$\leftrightarrow$FLIGHT} przesuwają granice faz o~1--2~klatki, co wpływa na czas kontaktu i~czas lotu, ale nie zmienia przypisania nogi.

Tabela~\ref{tab:error-breakdown} zestawia strukturę błędów wszystkich czterech modeli.

\begin{table}[htbp]
\centering
\caption{Struktura błędów klasyfikacji na zbiorze testowym}
\label{tab:error-breakdown}
\small
\begin{tabular}{lcccc}
\hline
\textbf{Model} & \textbf{Błędy} & \textbf{L$\leftrightarrow$R} & \textbf{S$\leftrightarrow$F} & \textbf{\% L$\leftrightarrow$R} \\
\hline
RF\,v1   & 556 & 153 & 403 & 27{,}5\% \\
RF\,v2   & 511 & 124 & 387 & 24{,}3\% \\
LSTM\,r1 & 422 &  75 & 347 & 17{,}8\% \\
LSTM\,r2 & 460 &  74 & 386 & 16{,}1\% \\
\hline
\end{tabular}
\end{table}

Kluczowa obserwacja: modele LSTM radykalnie redukują pomyłki L$\leftrightarrow$R --- z~153 (RF\,v1) do~75 (LSTM\,r1), czyli dwukrotnie. Jest to bezpośrednia konsekwencja kontekstu czasowego: LSTM widzi naprzemienną sekwencję faz (L$\to$F$\to$R$\to$F$\to$L) \\dzięki czemu `wie'', że po kontakcie lewej stopy następuje lot, a~potem kontakt prawej. Random Forest, klasyfikujący każdą klatkę niezależnie, nie ma tej informacji.

Inżynieria cech również pomaga: RF\,v2 popełnia mniej pomyłek L$\leftrightarrow$R niż RF\,v1 (124 vs 153), ponieważ normalizacja keypointów do układu ciała zmniejsza wpływ asymetrycznej perspektywy kamery na rozróżnienie stron.


\section{Ważność cech w~modelach Random Forest}
\label{sec:feature-importance}

W~tej sekcji sprawdzam, na czym właściwie modele Random Forest opierają swoje decyzje --- czyli które z~cech wejściowych okazują się dla nich najważniejsze. Pomaga to ocenić, czy model rzeczywiście patrzy na ruch biegacza, czy raczej zapamiętał, w~którym miejscu kadru zwykle pojawiają się stopy. Porównanie RF\,v1 i~RF\,v2 jest tu szczególnie ciekawe, bo drugi z~nich dostaje już przetworzone cechy (kąty stawów, znormalizowane pozycje keypointów względem ciała), a~pierwszy operuje na surowych współrzędnych z~MediaPipe.

Wagę każdej cechy Random Forest oblicza sprawdzając, jak skutecznie ta cecha pomaga drzewom rozdzielać próbki na właściwe klasy. Im częściej cecha sprawdzała się przy takich podziałach, tym wyższa jej waga. Wagi wszystkich cech sumują się do 100\%, więc wartość 3{,}6\% przy konkretnej cesze oznacza, że odpowiada ona za 3{,}6\% decyzji całego lasu. To pomocnicza miara --- mówi nam, na co model patrzył przy klasyfikacji, ale nie udowadnia, że właśnie ta cecha fizycznie decyduje o~tym, czy klatka to kontakt, czy lot.

Tabela poniżej zestawia pięć najważniejszych cech w~obu modelach.

\begin{table}[htbp]
\centering
\caption{Pięć najważniejszych cech w~modelach Random Forest}
\label{tab:feature-importance}
\small
\begin{tabular}{clcclc}
\hline
\multicolumn{3}{c}{\textbf{RF\,v1 (surowe keypointy)}} & \multicolumn{3}{c}{\textbf{RF\,v2 (cechy inżynierowane)}} \\
\ & Cecha & Waga & \ & Cecha & Waga \\
\hline
1 & RIGHT\_FOOT\_INDEX\_y & 3{,}6\% & 1 & left\_knee\_angle      & 3{,}6\% \\
2 & RIGHT\_ANKLE\_y       & 3{,}4\% & 2 & right\_knee\_angle     & 3{,}3\% \\
3 & RIGHT\_HEEL\_y        & 3{,}2\% & 3 & RIGHT\_HEEL\_y\_norm   & 3{,}2\% \\
4 & LEFT\_ANKLE\_y        & 3{,}1\% & 4 & RIGHT\_ANKLE\_y\_norm  & 2{,}8\% \\
5 & LEFT\_HEEL\_y         & 3{,}0\% & 5 & LEFT\_FOOT\_INDEX\_x\_norm & 2{,}5\% \\
\hline
\end{tabular}
\end{table}

W~RF\,v1 wszystkie pięć najważniejszych cech to surowe współrzędne $y$ stóp i~kostek. Model nauczył się, że gdy stopa jest nisko w~kadrze (czyli $y$ blisko 1 w~układzie MediaPipe), to prawdopodobnie dotyka ziemi. Jest to sensowne, ale działa dobrze tylko wtedy, gdy biegacz znajduje się w~podobnym miejscu kadru co w~filmach treningowych. Inne ujęcie, inna rozdzielczość, inny biegacz --- i wartości $y$ się przesuwają, a~model traci punkt odniesienia.

W~RF\,v2 na pierwszych dwóch miejscach znalazły się kąty kolan --- lewego i~prawego. Jest to dużo bardziej stabilna cecha niż surowe $y$. Kat kolana $170^\circ$ oznacza nogę prawię wyprostowaną niezależnie od tego, gdzie biegacz stoi w~kadrze i~jakiej jest postury. Na dalszych miejscach pojawiają się cechy z~dopiskiem \texttt{\_norm}, czyli pozycje keypointów przeliczone względem środka bioder i~długości tułowia --- również odporne na zmianę położenia w~kadrze.

To porównanie pokazuje, dlaczego RF\,v2 lepiej radzi sobie z~nowymi nagraniami. RF\,v1 zapamiętał wzorce wyglądu konkretnych biegaczy w~konkretnych ujęciach i~na nowych filmach jest zagubiony. RF\,v2 dostał cechy opisujące samą technikę biegu, a~nie jej konkretną wizualną realizację --- i~dlatego rozpoznaje tę technikę także u~innych biegaczy.



\section{Badanie wpływu korekty proporcji obrazu}
\label{sec:ablation-aspect}

Aby ocenić wpływ korekty proporcji obrazu (sekcja~\ref{sec:aspect-fix}), porównałem wyniki modelu LSTM\,r1 przed i~po zastosowaniu korekty. 

\begin{table}[htbp]
\centering
\caption{Wpływ korekty proporcji obrazu na dokładność LSTM\,r1}
\label{tab:aspect-fix-results}
\small
\begin{tabular}{lcccc}
\hline
 & \textbf{Bez korekty} & \textbf{Z~korektą} & \textbf{$\Delta$} & \textbf{Orientacja} \\
\hline
Film~2 (13~km/h)      & 54{,}9\% & 54{,}5\% & $-$0{,}4~pp & pozioma \\
Film~9 (zmienne tempo) & 68{,}1\% & 70{,}9\% & $+$2{,}8~pp & pozioma \\
Film~10 (kamera pionowa) & 75{,}8\% & 85{,}9\% & $+$10{,}1~pp & \textbf{pionowa} \\
\hline
\textbf{Łącznie (test)} & \textbf{67{,}1\%} & \textbf{70{,}9\%} & $+$\textbf{3{,}8~pp} & \\
\hline
\end{tabular}
\end{table}

Korekta daje największą poprawę na filmie~10 (wideo pionowe, 608$\times$1080~px): +10{,}1~pp. Wynika to z~faktu, że w~formacie pionowym proporcje szerokości do wysokości ($608/1080 \approx 0{,}56$) znacznie odbiegają od proporcji filmów poziomych ($\approx$1{,}33--1{,}78), co bez korekty powoduje błędne obliczanie długości tułowia i~kątów stawów. Dla filmów o~standardowych proporcjach (filmy~2 i~9) wpływ korekty jest niewielki ($-$0{,}4 do $+$2{,}8~pp).

Łączna poprawa na zbiorze testowym wyniosła +3{,}8~pp (z~67{,}1\% do~70{,}9\%). Korekta jest prostym krokiem preprocessingu --- przemnożeniem współrzędnych przez wymiary kadru ($x$ przez szerokość, $y$ przez wysokość) --- który nie wydłuża czasu inferencji, a~istotnie poprawia wyniki na nagraniach o~nietypowych proporcjach.

\section{Analiza wyników per-film}
\label{sec:per-film}

Tabela~\ref{tab:per-film} przedstawia dokładność wszystkich modeli na poszczególnych filmach testowych. Wyniki uwidaczniają, jak różnorodność warunków nagrania wpływa na jakość klasyfikacji.

\begin{table}[H]
\centering
\caption{Dokładność per-film na zbiorze testowym}
\label{tab:per-film}
\small
\begin{tabular}{lcccc}
\hline
\textbf{Film} & \textbf{RF\,v1} & \textbf{RF\,v2} & \textbf{LSTM\,r1} & \textbf{LSTM\,r2} \\
\hline
2 --- bieg 13\,km/h ($\approx$~300~klatek)      & 51{,}7\% & 60{,}3\% & 54{,}5\% & 54{,}9\% \\
9 --- zmienne tempo ($\approx$~870~klatek)      & 63{,}7\% & 63{,}2\% & 70{,}9\% & 69{,}3\% \\
10 --- kamera pionowa ($\approx$~320~klatek)    & 70{,}3\% & 77{,}5\% & \textbf{85{,}9\%} & 77{,}8\% \\
\hline
Łącznie                               & 62{,}7\% & 65{,}7\% & \textbf{70{,}9\%} & 68{,}2\% \\
\hline
\end{tabular}
\end{table}

\textbf{Film~2} (bieg 13~km/h, niska rozdzielczość 360$\times$450~px) okazał się najtrudniejszy dla wszystkich modeli --- najlepszy z~nich (RF\,v2) osiągnął na nim jedynie 60{,}3\%, a~główny model pracy (LSTM\,r1) --- 54{,}5\%. Prawdopodobną przyczyną jest niska rozdzielczość, która ogranicza precyzję detekcji keypointów przez MediaPipe, oraz fakt, że biegacz biega w~tempie i~stylu niereprezentowanym w~zbiorze treningowym.

\textbf{Film~9} (bieg zmiennym tempem, od~0{,}8 do~3{,}5~m/s) to najdłuższe nagranie w~zbiorze testowym ($\approx$~870~klatek), przez co najmocniej waży na wyniku łącznym --- dokładność LSTM\,r1 na tym filmie (70{,}9\%) niemal pokrywa się z~wynikiem ogólnym. Jest to przypadek pośredni między trudnym filmem~2 a~łatwym filmem~10. Zmienne tempo (od marszu po szybki bieg) sprawia, że czas trwania faz zmienia się w~obrębie nagrania, a~mimo to oba modele sekwencyjne wyraźnie przewyższają tu Random Forest (LSTM\,r1 70{,}9\% wobec około 63\% obu wariantów RF) --- kontekst czasowy pomaga zwłaszcza przy zmiennej dynamice biegu.

\textbf{Film~10} (kamera pionowa, analiza z~fizjoterapeutą) to film, na którym LSTM\,r1 osiągnął najlepszy wynik (85{,}9\%). Wysoka jakość nagrania (dobra rozdzielczość, stabilna kamera) i~wyraźny bieg w~stałym tempie sprzyjają zarówno detekcji pozy, jak i~klasyfikacji faz.

Różnica między filmami pokazuje, że jakość pipeline'u silnie zależy od warunków nagrania --- jest to ograniczenie omówione w~rozdziale~\ref{ch:wrazliwosc}.

\section{Wpływ dokładności klasyfikatora na precyzję współczynników}
\label{sec:wplyw-na-metryki}

Dokładność klasyfikatora faz wpływa bezpośrednio na precyzję obliczanych współczynników biomechanicznych. Każda błędna klasyfikacja klatki przesuwa granicę między fazami, co zmienia czas trwania segmentów STANCE i~FLIGHT.

Przy 30~FPS jedna klatka odpowiada $\approx$~33~ms. Jeśli model błędnie zaklasyfikuje 2~klatki na granicy STANCE$\to$FLIGHT, czas kontaktu (GCT) zostanie zawyżony o~$\approx$~67~ms --- przy typowym GCT wynoszącym 220~ms jest to błąd $\approx$~30\%. Przy 70{,}9\% dokładności LSTM\,r1 popełnia błędy STANCE$\leftrightarrow$FLIGHT w~$\approx$~347 z~$\approx$~1\,450~klatek, natomiast RF\,v1 --- w~$\approx$~403 z~$\approx$~1\,500. Poprawa o~5~pp oznacza w~praktyce mniej przesuniętych granic faz i~dokładniejsze czasy kontaktu.

Powyższy argument ma charakter teoretyczny, ale znalazł praktyczne potwierdzenie w~walidacji względem zegarka Garmin (sekcja~\ref{sec:walidacja-garmin}): na tym samym nagraniu kadencja (metryka oparta na zliczaniu kontaktów) zgodziła się z~pomiarem referencyjnym w~granicach 1\%, podczas gdy czas kontaktu (metryka oparta na pomiarze czasu trwania fazy) był zaniżony o~14\%. Potwierdza to, że to właśnie metryki czasu trwania --- zależne od dokładnego położenia granicy faz --- są najbardziej wrażliwe na błędy klasyfikatora.

Pomyłki L$\leftrightarrow$R mają jeszcze poważniejsze konsekwencje: jeśli model zamieni \texttt{LEFT\_STANCE} z~\texttt{RIGHT\_STANCE} w~kilku cyklach, wskaźnik symetrii Robinsona zostanie zaburzony --- system może fałszywie zasygnalizować asymetrię, której biegacz w~rzeczywistości nie ma. Dwukrotna redukcja pomyłek L$\leftrightarrow$R w~LSTM\,r1 (75) względem RF\,v1 (153) jest zatem kluczowa dla wiarygodności analizy symetrii.
\newpage
\section{Wnioski --- wybór modelu głównego}
\label{sec:wnioski-klasyfikator}

Na podstawie tych wyników jako model główny pipeline'u wybrałem \textbf{LSTM\,r1 z~korektą proporcji obrazu}, z~trzech powodów:

\begin{enumerate}
    \item \textbf{Najwyższa dokładność} testowa (70{,}9\%) i~F1 (0{,}709) spośród wszystkich modeli.
    \item \textbf{Najmniej pomyłek L$\leftrightarrow$R} (75; 17{,}8\% błędów), co przekłada się na wiarygodniejsze wskaźniki symetrii.
    \item \textbf{Najlepszy wynik na filmie~10} (85{,}9\%), potwierdzający skuteczność korekty proporcji przy nietypowej orientacji.
\end{enumerate}

Pozostałe modele zachowują wartość pomocniczą: RF\,v2 (65{,}7\%) jest szybki w~treningu i~nie wymaga GPU, więc nadaje się do prototypowania, a~LSTM\,r2 (68{,}2\%) --- choć nieco mniej dokładny --- jest mniejszy (187\,779 parametrów) i~oszczędniejszy obliczeniowo. Wspólnym ograniczeniem pozostaje niska dokładność na filmie~2 ($\approx$55\%), wskazująca na potrzebę rozszerzenia zbioru o~nagrania niskiej rozdzielczości i~zróżnicowanych stylów biegu.

% ============================================================================
% PLIK NIEAKTYWNY — rozdział 5 został usunięty z pracy decyzją usera (2026-05-28).
% Powód: 3 case studies bez statystyki nie stanowiły rzeczywistej walidacji,
% a poszczególne treści dublują się z innymi rozdziałami:
%   - Lista 12 współczynników + geneza → rozdz. 3.5
%   - Walidacja literaturowa → rozdz. 3.5 (wartości referencyjne)
%   - Case study Film 10 → rozdz. 4 (aspect ratio fix)
%   - Janek edge case → rozdz. 3.6 (reguła łączona SI+confidence)
%   - Ograniczenia obliczeń → rozdz. 8 (dyskusja)
%   - Pipeline pod różnymi warunkami → rozdz. 7 (wrażliwość)
%
% Plik zachowany jako draft do ewentualnej kanibalizacji fragmentów.
% Nie jest inputowany w main.tex.
% ============================================================================

\chapter{Współczynniki biegu --- wyniki i~walidacja}
\label{ch:wspolczynniki}

\section{Wprowadzenie}
\label{sec:wsp-wprowadzenie}

Niniejszy rozdział odpowiada na pierwsze pytanie badawcze postawione we wstępie: jak zaprojektować pipeline maszynowego uczenia, który na podstawie nagrania z~jednej kamery (\textit{monocular 2D}) wylicza zestaw współczynników biomechanicznych opisujących technikę biegu. W~rozdziale~\ref{ch:materialy-metody} (sekcja~\ref{sec:obliczanie-wspolczynnikow}) opisałem metody obliczania każdego współczynnika z~osobna. Tutaj pokazuję, jakie wartości pipeline produkuje w~praktyce, kiedy zostaje uruchomiony na rzeczywistych nagraniach, oraz oceniam ich wiarygodność przez porównanie z~zakresami raportowanymi w~literaturze biomechanicznej.

Pipeline wylicza łącznie 12~współczynników podzielonych na dwie grupy, wraz z~dodatkowymi wskaźnikami symetrii lewo-prawo.

\textbf{Pięć współczynników temporalnych} obliczanych z~sekwencji etykiet faz biegu i~częstotliwości klatek wideo:
\begin{itemize}
    \item kadencja,
    \item czas kontaktu z~podłożem (GCT) --- osobno dla lewej i~prawej nogi,
    \item czas lotu,
    \item czas cyklu (oraz długość kroku, jeśli użytkownik poda prędkość bieżni),
    \item duty factor.
\end{itemize}

\textbf{Siedem współczynników przestrzennych} obliczanych z~geometrii keypointów MediaPipe Pose:
\begin{itemize}
    \item trzy kąty stawów (kolano, biodro, kostka) --- na każdą nogę osobno, uśrednione w~obrębie fazy,
    \item kąt kolana w~momencie kontaktu z~podłożem (\textit{initial contact}),
    \item pochylenie tułowia,
    \item oscylacja pionowa bioder, znormalizowana długością tułowia,
    \item wzorzec lądowania stopy (\textit{heel} / \textit{midfoot} / \textit{forefoot strike}).
\end{itemize}

Dodatkowo dla każdego współczynnika obliczanego osobno dla lewej i~prawej strony ciała wyznaczam wskaźnik symetrii Robinsona (SI), opisany w~sekcji~\ref{sec:symmetry}.

Pierwotnie planowałem dużo węższy zakres --- cztery do pięciu podstawowych metryk (kadencja, GCT, długość kroku, wzorzec lądowania stopy). W~trakcie pracy nad pipeline okazało się jednak, że kolejne współczynniki dają się obliczyć w~naturalny sposób z~tych samych danych wejściowych (sekwencja faz + wygładzone keypointy), bez dodatkowych kosztów obliczeniowych ani dodatkowych wymagań wobec nagrania. Rozszerzony zestaw 12~współczynników daje pełniejszy obraz techniki biegu i~pozwala silnikowi reguł rekomendacji (rozdział~\ref{ch:rekomendacje}) wykrywać szerszy zakres potencjalnych problemów.

\section{Strategia walidacji}
\label{sec:strategia-walidacji}

Najbardziej miarodajną metodą walidacji byłoby porównanie wartości wyliczonych z~monocular 2D wideo z~pomiarami trójwymiarowego systemu motion capture (3D MoCap), takiego jak Vicon lub Qualisys, traktowanymi jako ground truth. Dostęp do takiego laboratorium wymaga znacznego budżetu i~specjalistycznej infrastruktury, co wykracza poza możliwości pracy magisterskiej realizowanej w~warunkach domowych. Z~tego samego powodu w~zbiorze nie ma żadnego nagrania, dla którego istniałyby równoległe pomiary 3D, więc nie mogę raportować precyzji w~jednostkach fizycznych (np.~MAE w~milisekundach lub stopniach).

Zastosowałem walidację pośrednią: dla każdego współczynnika sprawdzam, czy wyliczona wartość mieści się w~zakresie publikowanym w~literaturze biomechanicznej --- przede wszystkim Novacheck~\cite{novacheck1998}, Heiderscheit~\cite{heiderscheit2011}, Souza~\cite{souza2016} i~Daoud~\cite{daoud2012}. Zakresy te są zebrane w~tabeli referencyjnej (sekcja~\ref{sec:obliczanie-wspolczynnikow}) i~służą zarówno jako kryterium oceny, jak i~jako progi dla silnika reguł rekomendacji.

Walidację jakościową przeprowadziłem na trzech biegaczach o~świadomie zróżnicowanej technice: zaawansowanym (Adam, zbiór treningowy), pośrednim (film~10 ze~zbioru testowego) i~z~wyraźnymi problemami (Janek, przypadek brzegowy spoza zbioru). Celem nie było dowodzenie pełnej precyzji pipeline'u, lecz sprawdzenie, czy potrafi on \textit{różnicować} biegaczy --- czyli czy lepsza technika prowadzi do wartości bliższych zakresom z~literatury, a~problemy techniczne wywołują sygnały ostrzegawcze.

\section{Studium przypadków --- trzech biegaczy}
\label{sec:case-studies}

Tabela zbiorcza wszystkich wartości znajduje się w~sekcji~\ref{sec:porownanie-zbiorcze}. Tutaj przedstawiam każdy przypadek osobno, koncentrując się na wnioskach dotyczących pipeline'u, a~nie na pełnej charakterystyce techniki biegu konkretnego biegacza.

\subsection{Adam --- biegacz zaawansowany}
\label{sec:case-adam}

Pierwszym biegaczem walidacyjnym jest Adam, członek klubu biegowego Zabiegani na Szamę. Jego nagranie (film~12, 80~s, FPS~$\approx$~30) trafiło do zbioru treningowego, więc dla wyników klasyfikatora faz nie traktuję go jako oceny generalizacji. W~kontekście walidacji współczynników biomechanicznych jest jednak przydatne jako sprawdzian, czy pipeline produkuje sensowne wartości dla biegacza zaawansowanego.

Wszystkie wskaźniki temporalne mieszczą się w~zakresach typowych dla biegacza zaawansowanego: kadencja 173~spm (literatura: 170--185), GCT lewej nogi 203~ms i~prawej 245~ms (zakres rekreacyjnego i~szybkiego biegu, 160--280~ms), oscylacja pionowa znormalizowana 0{,}14 długości tułowia (dobry biegacz: 0{,}12--0{,}16). Foot strike na obu nogach to forefoot, z~wyraźną dominacją w~96\% kontaktów lewej i~75\% prawej. Kąt kolana w~momencie kontaktu (161° lewa, 162° prawa) mieści się w~zalecanym zakresie 160--175°.

Jedna obserwacja jest jednak nietrywialna: wskaźnik symetrii GCT wynosi SI~=~18{,}7\%, co przekracza próg 10\% i~zostało oznaczone jako potencjalny problem. Pipeline samodzielnie nie rozstrzyga, czy biegacz rzeczywiście ma asymetryczne kontakty, czy jest to artefakt klasyfikatora przesuwającego granice faz o~1--2~klatki. Analiza wrażliwości (rozdział~\ref{ch:wrazliwosc}) pokazuje, że w~tym wypadku różnica wynika w~znacznej mierze z~drugiego źródła. Przykład ten dobrze ilustruje, że walidacja przez zakresy literatury mówi nam, \textit{kiedy} wartość wykracza poza normę, ale nie zawsze \textit{dlaczego}.

\subsection{Film 10 --- analiza z~fizjoterapeutą}
\label{sec:case-film10}

Drugim biegaczem walidacyjnym jest osoba z~filmu~10 (zbiór testowy, $\approx$~320~klatek, FPS~$\approx$~30, ujęcie pionowe $608 \times 1080$~px). Wybrałem ten przypadek celowo, ponieważ orientacja kamery uwypukla wrażliwość różnych grup współczynników na warunki akwizycji.

Wartości temporalne pozostają sensowne: kadencja 163~spm (zakres rekreacyjny), GCT 285/226~ms (jogging do szybkiego biegu), czas lotu 116~ms, pochylenie tułowia 9{,}9°. Pipeline temporalny radzi sobie z~filmem pionowym, ponieważ czasy faz nie zależą od orientacji kadru.

Współczynniki przestrzenne pokazują jednak limity podejścia 2D. Kąt kolana w~momencie kontaktu wynosi 98° (lewa) i~137° (prawa) --- obie wartości daleko poza zakresem prawidłowym 160--175° i~oznaczone jako problem ``siedzącego biegu''. Realnie wynika to z~połączenia perspektywy kamery z~techniką biegacza: w~ujęciu pionowym geometria kątów na płaszczyźnie obrazu jest zniekształcona względem rzeczywistej pozycji nogi w~przestrzeni. Najbardziej spektakularne są kąty stopy w~detekcji wzorca lądowania --- $-97^\circ$ i~$-98^\circ$ w~wartości bezwzględnej, przekraczające próg wiarygodności $45^\circ$ wielokrotnie. Pipeline poprawnie zapala flagę \texttt{low\_confidence} (sekcja~\ref{sec:obliczanie-wspolczynnikow}) i~oznacza wzorzec lądowania jako niewiarygodny. Ten przypadek pokazuje istotną cechę projektową: pipeline jest świadomy własnych ograniczeń i~nie udaje precyzji, której fizycznie nie może zapewnić.

\subsection{Janek --- przypadek brzegowy}
\label{sec:case-janek}

Trzeci biegacz, Janek, jest przypadkiem szczególnym. Jego nagranie (FPS~30, $\approx$~2400~klatek) nie znalazło się w~żadnym ze~zbiorów --- ani treningowym, ani walidacyjnym, ani testowym --- ponieważ posłużyło wyłącznie do oceny pipeline'u jako całości na nieznanym wcześniej materiale.

Wyniki pipeline'u są wyraźnie odmienne od dwóch poprzednich przypadków: kadencja 148~spm (poniżej progu 150 oznaczonego jako bardzo niska, sygnalizująca overstriding lub chód), GCT lewej nogi 192~ms wobec prawej 357~ms (różnica 165~ms), duty factor lewa 0{,}24 wobec prawej 0{,}44, kąt kolana w~kontakcie 138° lewa i~148° prawa (oba poniżej zakresu 160--175°). Wskaźnik symetrii Robinsona dla GCT wynosi SI~=~59{,}8\%, dla duty factor 60{,}0\%. Wartości tego rzędu nie są spotykane u~zdrowych biegaczy --- to różnice z~obszaru patologii ruchu albo poważnego błędu pomiaru.

Pipeline poprawnie zareagował: silnik reguł rekomendacji wygenerował flagę \texttt{critical} dla asymetrii kroków oraz dla łącznej kombinacji wysokiego SI z~obniżoną pewnością predykcji (\texttt{avg\_confidence}~=~0{,}88), korzystając z~reguły łączonej opisanej w~sekcji~\ref{sec:rules-combined}. Ta sygnalizacja jest tu właściwa niezależnie od tego, czy źródłem są błędy klasyfikatora na nietypowym nagraniu, czy realny problem z~techniką biegu --- w~obu przypadkach wynik wymaga interwencji człowieka i~nie powinien być prezentowany użytkownikowi jako wiarygodna analiza biomechaniczna.

\section{Porównanie zbiorcze}
\label{sec:porownanie-zbiorcze}

Tabela~\ref{tab:wsp-zbiorcza} zestawia kluczowe współczynniki dla trzech biegaczy walidacyjnych.

\begin{table}[htbp]
\centering
\caption{Wyniki obliczeń współczynników dla trzech biegaczy walidacyjnych}
\label{tab:wsp-zbiorcza}
\small
\begin{tabular}{lccc}
\hline
\textbf{Współczynnik} & \textbf{Adam} & \textbf{Film~10} & \textbf{Janek} \\
\hline
Kadencja [spm]                       & 173      & 163      & 148 (🔴) \\
GCT lewa [ms]                        & 203      & 285      & 192 \\
GCT prawa [ms]                       & 245      & 226      & 357 (🔴) \\
Czas lotu [ms]                       & 122      & 116      & 130 \\
Duty factor L / P                    & 0{,}30 / 0{,}36 & 0{,}39 / 0{,}31 & 0{,}24 / 0{,}44 \\
Kąt kolana @ kontakt L [$^\circ$]    & 162      & 98 (🟡)  & 138 (🟡) \\
Kąt kolana @ kontakt P [$^\circ$]    & 162      & 137 (🟡) & 149 (🟡) \\
Pochylenie tułowia [$^\circ$]        & 2{,}3 (🟡) & 9{,}9    & 9{,}9 \\
Vertical osc. (per torso)            & 0{,}14   & 0{,}10 (🟡) & 0{,}11 (🟡) \\
Foot strike L / P                    & FF / FF  & FF / FF$^\dagger$ & MF / MF \\
SI GCT [\%]                          & 18{,}7   & 22{,}8   & 59{,}9 (🔴) \\
SI duty factor [\%]                  & 18{,}7   & 21{,}4   & 60{,}0 (🔴) \\
Flagi krytyczne (\texttt{critical})  & 0        & 0        & 2 \\
\hline
\multicolumn{4}{l}{\footnotesize FF = forefoot, MF = midfoot. $^\dagger$ Flaga \texttt{low\_confidence} (kąt stopy $>$~$45^\circ$ b.w.).} \\
\end{tabular}
\end{table}

Z~zestawienia wynikają trzy obserwacje. Po pierwsze, pipeline różnicuje trzech biegaczy: każda kolumna ma inną sygnaturę odchyleń od literatury i~inną liczbę zapalonych flag krytycznych. Po drugie, wszystkie wartości temporalne mieszczą się w~zakresach raportowanych w~literaturze i~odpowiadają intuicji biomechanicznej --- Adam ma najwyższą kadencję i~najkrótsze kontakty, Janek najniższą kadencję i~najbardziej asymetryczne kontakty. Po trzecie, współczynniki przestrzenne (zwłaszcza kąty kolana w~momencie kontaktu i~wzorzec lądowania stopy) okazują się wyraźnie wrażliwe na warunki nagrania --- u~biegacza z~filmu~10 dają wartości fizycznie nieprawdopodobne, ale pipeline rozpoznaje to przez flagę \texttt{low\_confidence}.

Wskaźnik symetrii Robinsona poprawnie wskazuje skalę problemu: 18--22\% u~biegaczy bez większych zastrzeżeń biomechanicznych i~60\% w~przypadku Janka, gdzie ani biegacz, ani warunki nagrania nie są typowe.

\section{Ograniczenia obliczeń}
\label{sec:ograniczenia-obliczen}

\begin{itemize}
    \item \textbf{Brak kalibracji piksel--metr.} Oscylacja pionowa raportowana jest w~jednostkach znormalizowanych długością tułowia, a~nie w~centymetrach. Przeliczenie na centymetry wymagałoby znajomości wzrostu biegacza i~odległości kamery od niego, co rzadko bywa dostępne w~praktycznym scenariuszu użycia.
    \item \textbf{Monocular 2D --- artefakty perspektywy.} Wszystkie współczynniki kątowe wyliczane są w~płaszczyźnie obrazu, nie w~przestrzeni trójwymiarowej. Przy ujęciu innym niż ściśle boczne (np.~wideo pionowe z~filmu~10) geometria kątów ulega zniekształceniu. Pipeline kompensuje to częściowo przez flagę \texttt{low\_confidence} dla kąta stopy, ale nie dla kątów stawów.
    \item \textbf{Stride length wymaga inputu użytkownika.} Na bieżni mechanicznej biegacz porusza się w~miejscu, więc długości kroku nie da się wyliczyć z~samego wideo --- pipeline wymaga podania prędkości bieżni przez użytkownika. Dla nagrań bez tej informacji długość kroku nie jest raportowana.
    \item \textbf{Brak ground truth dla precyzji.} Walidacja przez zakresy literatury pozwala stwierdzić, czy wartość jest realistyczna biomechanicznie, ale nie pozwala podać dokładności rzędu MAE w~jednostkach fizycznych. Dla pełnej oceny wymagana byłaby równoległa rejestracja 3D MoCap, co pozostawiam jako kierunek dalszych prac.
    \item \textbf{Non-determinizm MediaPipe.} Domyślny silnik (XNNPACK delegate) wprowadza drobne, nieprzewidywalne różnice między uruchomieniami dla tego samego nagrania. Przy 30~FPS i~$\approx$~12~klatkach na cykl różnice te są zwykle pomijalne, ale w~filmach z~niskim FPS mogą się kumulować.
\end{itemize}

\chapter{System rekomendacji biomechanicznych}
\label{ch:rekomendacje}

Pipeline kończy się modułem, który przekształca obliczone współczynniki w~konkretne wskazówki treningowe. Rozdział ten opisuje, dlaczego rekomendacje oparłem na regułach z~literatury zamiast na uczeniu maszynowym, jak zbudowany jest silnik reguł, jakie reguły zawiera oraz jak wygląda wygenerowany raport dla rzeczywistego biegacza.

\section{Dlaczego reguły, a~nie uczenie maszynowe}
\label{sec:rule-based-vs-ml}

Rekomendacje nie są uczone z~danych --- są zakodowanymi na stałe regułami odzwierciedlającymi zakresy referencyjne i~wzorce ryzyka ustalone w~literaturze biomechanicznej. Wybór ten wynika z~czterech przesłanek:

\begin{itemize}
    \item \textbf{Rozmiar zbioru danych.} Zbiór 15~nagrań i~$\approx$~12~biegaczy jest o~rzędy wielkości za mały, aby wytrenować model przewidujący rekomendacje. Model uczony na tak małej próbce nauczyłby się raczej idiosynkrazji konkretnych biegaczy niż ogólnych zasad techniki biegu.
    \item \textbf{Wymóg źródła.} Rekomendacja dotycząca techniki biegu funkcjonuje w~kontekście zbliżonym do porady quasi-medycznej. Każda wskazówka powinna mieć możliwe do wskazania źródło w~literaturze --- użytkownik (lub współpracujący z~nim trener czy fizjoterapeuta) musi móc zweryfikować, skąd wzięła się dana sugestia. Model uczony maszynowo takiego uzasadnienia nie dostarcza.
    \item \textbf{Interpretowalność.} W~podejściu regułowym każda rekomendacja ma jawnie podany próg, zmierzoną wartość i~publikację, z~której pochodzi. Nie ma ``czarnej skrzynki'' --- ścieżkę od pomiaru do wskazówki da się prześledzić krok po kroku.
    \item \textbf{Transparentność i~powtarzalność.} Reguły są deterministyczne: te same wartości wejściowe zawsze dają tę samą rekomendację. Ułatwia to walidację, debugowanie i~rozszerzanie systemu o~kolejne reguły bez ryzyka nieprzewidzianych interakcji.
\end{itemize}

Literatura biomechaniczna dostarcza przy tym dobrze udokumentowane i~wielokrotnie potwierdzone zakresy referencyjne dla większości interesujących mnie współczynników, więc rezygnacja z~uczenia maszynowego nie oznacza utraty wiarygodności.

\section{Architektura silnika reguł}
\label{sec:rules-arch}

Silnik składa się z~11~funkcji reguł, z~których każda analizuje jeden aspekt techniki biegu (np.~kadencję, czas kontaktu, pochylenie tułowia). Każda funkcja otrzymuje na wejściu wyniki obliczeń współczynników i~zwraca listę rekomendacji --- obiektów zawierających:

\begin{itemize}
    \item \textbf{Poziom ważności} --- czterostopniowa skala:
    \begin{itemize}
        \item \texttt{critical} --- wartość wymagająca interwencji. Flaga ta zapala się w~dwóch sytuacjach: po pierwsze, gdy wynik wskazuje na rzeczywiste ryzyko biomechaniczne (np.~kąt kolana powyżej $175^\circ$ przy kontakcie = overstriding), po drugie, gdy wynik sugeruje, że klasyfikator faz coś źle sklasyfikował i~obliczone współczynniki są niewiarygodne (np.~asymetria liczby kroków lewej i~prawej nogi przekraczająca~20\%, co świadczy o~tym, że model myli fazy). Oba przypadki wymagają reakcji człowieka, ale innego rodzaju: w~pierwszym biegacz powinien skorygować technikę, w~drugim --- należy przenagrać film z~lepszą perspektywą kamery lub w~lepszych warunkach,
        \item \texttt{warning} --- wartość poza zalecanym zakresem, ale nie krytyczna,
        \item \texttt{watch} --- wartość w~zakresie akceptowalnym, ale warta monitorowania,
        \item \texttt{info} --- informacja pozytywna lub neutralna (np.~``kadencja w~zalecanym zakresie''),
    \end{itemize}
    \item \textbf{Źródło literaturowe}  --- autor i~rok publikacji, na której oparta jest reguła,
    \item \textbf{Szczegóły pomiaru} --- konkretne zmierzone wartości i~zakresy referencyjne,
    \item \textbf{Sugestia} --- konkretne działanie korygujące (np.~``zwiększ kadencję o~5--10\%'').
\end{itemize}

Pod względem implementacyjnym silnik jest zwykłym modułem Pythona: funkcje reguł zwracają obiekty rekomendacji, a~osobna funkcja składa je w~raport w~formacie Markdown. Moduł jest wywoływany jako ostatni krok pipeline'u analizy, po obliczeniu wszystkich współczynników.

\section{Katalog reguł}
\label{sec:rules-categories}

Tabela~\ref{tab:rules} przedstawia reguły z~przykładowymi progami i~źródłami literaturowymi, pogrupowane na trzy rodzaje: reguły temporalne (operujące na czasach faz), przestrzenne (operujące na geometrii keypointów) oraz walidacyjne (oceniające wiarygodność samej predykcji).

\begin{table}[htbp]
\centering
\caption{Katalog reguł rekomendacji z~przykładowymi progami}
\label{tab:rules}
\small
\begin{tabular}{p{3.2cm}p{5.5cm}p{3.8cm}}
\hline
\textbf{Aspekt} & \textbf{Przykładowa reguła (próg $\to$ poziom)} & \textbf{Źródło} \\
\hline
\multicolumn{3}{l}{\textbf{Reguły temporalne} --- na podstawie czasu trwania faz biegu} \\
Kadencja           & $<$160~spm $\to$ warning (overstriding) & Heiderscheit 2011 \\
GCT                & $>$280~ms $\to$ warning (długi kontakt) & Souza 2016 \\
Duty factor        & $\geq$0{,}5 $\to$ critical (chód)       & Souza 2016 \\
Stride length      & $>$2{,}2~m $\wedge$ kadencja~$<$160 $\to$ overstriding & Heiderscheit 2011 \\
\hline
\multicolumn{3}{l}{\textbf{Reguły przestrzenne} --- na podstawie geometrii ciała w~kadrze} \\
Kolano przy kontakcie & $>175^\circ$ $\to$ critical (overstriding)   & Heiderscheit 2011 \\
Pochylenie tułowia & $<5^\circ$ $\to$ warning (zbyt pionowy)      & Novacheck 1998 \\
Oscylacja pionowa  & $>$0{,}24 torso $\to$ warning           & Novacheck 1998 \\
Symetria L/P       & SI~$>$10\% $\to$ warning                & Robinson 1987 \\
Foot strike        & L$\neq$R $\to$ warning (niespójność)    & Daoud 2012 \\
\hline
\multicolumn{3}{l}{\textbf{Reguły walidacyjne} --- wiarygodność samej predykcji} \\
Pewność predykcji  & avg\_confidence $<$0{,}85 $\to$ warning & walidacja wewnętrzna \\
Asymetria kroków   & $|$L$-$R$|$ / total $>$20\% $\to$ critical & walidacja wewnętrzna \\
\hline
\end{tabular}
\end{table}

\section{Reguły łączone}
\label{sec:rules-combined}

Część reguł analizuje kombinację dwóch współczynników jednocześnie, co pozwala wykryć wzorce niewidoczne przy analizie każdego współczynnika osobno. Zaimplementowałem trzy reguły łączone:

\begin{enumerate}
    \item \textbf{Overstriding (kadencja + GCT)}: kadencja poniżej~160~spm w~połączeniu ze~średnim GCT powyżej~280~ms to klasyczny wzorzec overstridingu --- stopa ląduje zbyt daleko przed środkiem ciężkości \cite{heiderscheit2011}.
    \item \textbf{Overstriding (stride + kadencja)}: stride powyżej~2{,}2~m przy kadencji poniżej~160~spm --- długi krok wynikający z~niskiej kadencji, a~nie z~dużej prędkości.
    \item \textbf{Jakość predykcji (SI + confidence)}: maksymalny SI powyżej~50\% w~połączeniu ze~średnią pewnością predykcji poniżej~0{,}90 --- sygnał, że asymetria L/P jest artefaktem klasyfikatora, nie rzeczywistą cechą biegacza. Reguła ta powstała na podstawie analizy przypadku brzegowego jednego z~biegaczy, na których walidowałem cały pipeline (poza zbiorem treningowym, walidacyjnym i~testowym modelu), u~którego żaden pojedynczy wskaźnik z~osobna nie przekraczał swojego progu ostrzegawczego, podczas gdy ich kombinacja --- wysoka asymetria GCT (60\%) przy obniżonej pewności predykcji --- ujawniała artefakt.
\end{enumerate}

Reguły łączone powstały z~obserwacji, że pojedyncze wskaźniki bywają zawodne na nietypowych nagraniach. Sprawdzenie dwóch warunków naraz wychwytuje sytuacje, które każdy z~nich osobno by nie wylapał.

\section{Źródła reguł}
\label{sec:rules-sources}

Zakresy referencyjne i~progi reguł oparłem na pięciu źródłach literaturowych oraz na walidacji wewnętrznej:
\begin{itemize}
    \item \textbf{Heiderscheit i~wsp.\ (2011)} \cite{heiderscheit2011} --- zwiększenie kadencji o~5--10\% obniża obciążenia stawów, kąt kolana przy kontakcie i~overstriding,
    \item \textbf{Novacheck (1998)} \cite{novacheck1998} --- pochylenie tułowia, kąty stawów, oscylacja pionowa, fazy cyklu biegowego,
    \item \textbf{Souza (2016)} \cite{souza2016} --- GCT, duty factor, foot strike w~kontekście analizy wideo,
    \item \textbf{Daoud i~wsp.\ (2012)} \cite{daoud2012} --- foot strike pattern a~częstość kontuzji,
    \item \textbf{Robinson i~wsp.\ (1987)} \cite{robinson1987} --- wskaźnik symetrii SI z~progami 5\%/10\%,
    \item \textbf{walidacja wewnętrzna} --- progi jakości predykcji (pewność predykcji, asymetria kroków) kalibrowane na biegaczach użytych do walidacji całego pipeline'u.
\end{itemize}
\newpage
\section{Przykład wygenerowanego raportu}
\label{sec:rules-example}

Aby pokazać, co system produkuje w~praktyce, przedstawiam raport wygenerowany dla biegacza Adama (film~12, kadencja 173~spm, pewność predykcji 0{,}91). Silnik zwrócił 8~rekomendacji: 0~krytycznych, 2~ostrzeżenia, 1~do monitorowania i~5~informacji pozytywnych (tabela~\ref{tab:rec-example}).

\begin{table}[H]
\centering
\caption{Rekomendacje wygenerowane dla biegacza Adama}
\label{tab:rec-example}
\small
\begin{tabular}{p{1.6cm}p{4.3cm}p{4.2cm}p{2.4cm}}
\hline
\textbf{Poziom} & \textbf{Rekomendacja} & \textbf{Pomiar} & \textbf{Źródło} \\
\hline
warning & Tułów zbyt pionowy & pochylenie $2{,}3^\circ$ (cel $\geq 5^\circ$) & Novacheck 1998 \\
warning & Asymetria L/P $>$10\% & GCT SI~=~18{,}7\%, duty factor SI~=~18{,}7\% & Robinson 1987 \\
watch   & Łagodna asymetria L/P & kąt kostki SI~=~5{,}6\% & Robinson 1987 \\
info    & Kadencja w~zakresie & 173~spm & Heiderscheit 2011 \\
info    & Kolano L prawidłowe & $161{,}8^\circ$ & Heiderscheit 2011 \\
info    & Kolano P prawidłowe & $162{,}0^\circ$ & Heiderscheit 2011 \\
info    & Oscylacja ekonomiczna & 0{,}140 torso & Novacheck 1998 \\
info    & Spójny foot strike & L~=~P~=~forefoot & Daoud 2012 \\
\hline
\end{tabular}
\end{table}

Każda rekomendacja zawiera nie tylko etykietę, ale też zmierzoną wartość, zakres referencyjny, źródło literaturowe oraz --- w~przypadku ostrzeżeń --- konkretną sugestię korygującą. Warto zwrócić uwagę, że ostrzeżenie o~asymetrii (SI~=~18{,}7\%) opatrzone jest w~pełnym raporcie zastrzeżeniem, że część asymetrii w~danych z~pojedynczej kamery może być artefaktem perspektywy --- system nie przedstawia więc swoich wyników jako pewnej diagnozy, lecz jako sygnał do dalszej obserwacji lub weryfikacji człowieka.

\newpage
\section{Ograniczenia systemu}
\label{sec:ograniczenia-rekomendacji}

\begin{itemize}
    \item \textbf{Progi ogólne, nie zindywidualizowane.} Zakresy referencyjne pochodzą z~literatury i~mają charakter populacyjny. Indywidualna sytuacja biegacza (historia kontuzji, budowa ciała, cel treningowy) może wymagać innych progów, których system nie uwzględnia.
    \item \textbf{Operowanie na wartościach średnich.} Reguły działają na uśrednionych współczynnikach z~całego nagrania, nie na poszczególnych cyklach. Chwilowe odchylenia (np.zmęczenie pod koniec biegu) mogą zostać uśrednione i~niezauważone.
    \item \textbf{Zależność od jakości współczynników.} Rekomendacje są tak dobre, jak dane wejściowe. Błędy klasyfikatora faz lub artefakty perspektywy propagują się do rekomendacji --- stąd reguły walidacyjne, które mają wychwytywać najbardziej rażące przypadki.
    \item \textbf{Konieczność aktualizacji przy rozszerzaniu.} Dodanie nowej reguły wymaga znalezienia i~zweryfikowania źródła literaturowego dla jej progów, co jest pracą manualną i~nie skaluje się automatycznie wraz z~danymi.
\end{itemize}

\chapter{Wrażliwość współczynników na warunki akwizycji wideo}
\label{ch:wrazliwosc}

Pipeline opisany w~poprzednich rozdziałach produkuje wartości współczynników niezależnie od tego, w~jakich warunkach nagrano wideo. Nie wszystkie współczynniki są jednak równie odporne na te warunki: część zależy od perspektywy kamery, część od częstotliwości klatek, a~niektóre pozostają wiarygodne niezależnie od obu. Rozdział ten bada tę wrażliwość i~pokazuje, że najbardziej zawodne przypadki da się wykryć automatycznie, bez dostępu do pomiarów referencyjnych.

\section{Hipotezy}
\label{sec:hipotezy-p3}

Rozdział weryfikuje trzy hipotezy dotyczące wpływu warunków akwizycji na jakość obliczanych współczynników:

\begin{enumerate}
    \item \textbf{Współczynniki przestrzenne są wrażliwe na perspektywę kamery.} Wartości oparte na geometrii keypointów w~płaszczyźnie obrazu (zwłaszcza wzorzec lądowania stopy) zniekształcają się przy ujęciu innym niż ściśle boczne.
    \item \textbf{Współczynniki temporalne są wrażliwe na częstotliwość klatek (FPS).} Przy niskim FPS kwantyzacja czasu na dyskretne klatki wprowadza istotny błąd względny do metryk czasowych.
    \item \textbf{Niską wiarygodność pomiaru można wykryć automatycznie, bez pomiarów referencyjnych} --- analizując, czy wyliczona wartość mieści się w~zakresie fizjologicznie możliwym.
\end{enumerate}

Każdą hipotezę badam na studium przypadku opartym na konkretnych nagraniach ze~zbioru.

\section{Studium 1: wzorzec lądowania stopy a~perspektywa kamery}
\label{sec:case1-foot-strike}

Wzorzec lądowania stopy (sekcja~\ref{sec:obliczanie-wspolczynnikow}) klasyfikowany jest na podstawie kąta wektora od~pięty do~czubka stopy w~klatce wejścia w~fazę kontaktu. Aby ocenić jego wiarygodność, przeprowadziłem walidację wizualną na trzech biegaczach o~różnych warunkach nagrania: Adamie (film~12, zbiór treningowy), biegaczu z~filmu~10 (zbiór testowy) oraz Janku (nagranie spoza zbiorów modelu). Ponieważ jest to walidacja wizualna samej geometrii kąta stopy, a~nie ocena generalizacji klasyfikatora, użycie filmu ze~zbioru treningowego (Adam) nie wpływa na wnioski. Dla każdego biegacza wyrenderowałem szkielet MediaPipe na klatkach wejścia w~fazę STANCE (po trzy klatki na nogę, rozłożone równomiernie) i~porównałem obliczony kąt z~rzeczywistym ułożeniem stopy widocznym na obrazie. Tabela~\ref{tab:foot-strike-validation} zestawia wyniki.

\begin{table}[htbp]
\centering
\caption{Walidacja wzorca lądowania stopy dla trzech warunków nagrania}
\label{tab:foot-strike-validation}
\small
\begin{tabular}{p{2.2cm}p{4.2cm}p{2.6cm}p{3.6cm}}
\hline
\textbf{Biegacz} & \textbf{Ujęcie} & \textbf{Średni kąt L / P} & \textbf{Werdykt} \\
\hline
Janek    & boczne, kamera prostopadle do toru & $-4^\circ$ / $-3^\circ$ & midfoot, wiarygodne \\
Adam     & boczne, kamera lekko od dołu & $-33^\circ$ / $-12^\circ$ & forefoot; lewa noga zawyżona przez perspektywę \\
Film~10  & pionowe (portrait) & $-97^\circ$ / $-99^\circ$ & bezwartościowe ($|\text{kąt}| > 90^\circ$) \\
\hline
\end{tabular}
\end{table}

Wyniki układają się w~czytelny wzorzec. Przy standardowym ujęciu bocznym (Janek) obliczone kąty zgadzają się z~obserwacją wizualną --- stopa ląduje płasko, a~klasyfikacja midfoot jest wiarygodna. Przy ujęciu bocznym, ale z~kamerą ustawioną lekko od dołu (Adam), prawa noga (bliżej kamery) daje sensowny wynik, natomiast lewa (dalej, silniej dotknięta perspektywą) ma kąt zawyżony --- wektor pięta--czubek jest skrócony w~rzucie, co matematycznie wzmacnia składową pionową. Asymetria $-33^\circ$ wobec $-12^\circ$ jest więc artefaktem perspektywy, nie biomechaniki. Przy ujęciu pionowym (Film~10) wszystkie kąty przekraczają $90^\circ$ w~wartości bezwzględnej, co jest geometrycznie niemożliwe dla biegacza zwróconego w~kierunku biegu --- predykcja jest bezwartościowa.

Studium to potwierdza pierwszą hipotezę: wzorzec lądowania stopy jest wiarygodny przy ujęciu ściśle bocznym i~zawodzi przy ujęciu pionowym lub ukośnym. Warto zauważyć, że ta sama pionowa orientacja Filmu~10, która psuje kąt stopy, obniża również dokładność klasyfikatora faz --- problem ten zaadresowałem korektą proporcji obrazu, której wpływ przeanalizowałem osobno w~sekcji~\ref{sec:ablation-aspect}.

\section{Studium 2: częstotliwość klatek a~precyzja metryk temporalnych}
\label{sec:case2-fps}

Współczynniki temporalne wyliczam z~liczby klatek przypadających na daną fazę (sekcja~\ref{sec:segmentacja}). Każda faza trwa całkowitą liczbę klatek, więc jej czas jest skwantowany do wielokrotności $1/\text{FPS}$. Przy wysokim FPS kwantyzacja jest drobna, ale przy niskim staje się istotnym źródłem błędu.

Rozważmy typowy czas kontaktu z~podłożem GCT~$=250$~ms. Przy 30~FPS jedna klatka odpowiada $\approx$~33~ms, więc kontakt obejmuje około 7,5~klatki, a~błąd $\pm 1$~klatki to $\pm 13\%$. Przy 13,33~FPS (jak w~slow-motion Filmie~7) jedna klatka to $\approx$~75~ms, kontakt obejmuje zaledwie $\approx$~3,3~klatki, a~ten sam błąd $\pm 1$~klatki rośnie do $\pm 30\%$. Im krótsza faza i~niższy FPS, tym większy względny błąd kwantyzacji.

Wniosek ma istotną konsekwencję metodologiczną, spójną z~założeniami przyjętymi przy budowie zbioru (sekcja~\ref{sec:dataset}): nagrania o~niskim FPS są pełnoprawnym materiałem \textit{treningowym} dla klasyfikatora faz (model uczy się pozycji ciała, nie czasu), ale nie nadają się do precyzyjnego wyliczania metryk \textit{temporalnych} na etapie inferencji. FPS należy więc traktować jako parametr jakości nagrania przy obliczaniu współczynników czasowych, a~nie przy treningu.

\section{Synteza: wrażliwość metryk i~możliwość auto-detekcji}
\label{sec:synteza-wrazliwosc}

Tabela~\ref{tab:wrazliwosc-synteza} zestawia wszystkie współczynniki pod kątem wrażliwości na orientację kamery i~na FPS oraz tego, czy daną wadliwość da się wykryć automatycznie.

\begin{table}[htbp]
\centering
\caption{Wrażliwość współczynników na warunki akwizycji i~możliwość automatycznej detekcji}
\label{tab:wrazliwosc-synteza}
\small
\begin{tabular}{p{3.4cm}p{2.8cm}p{2.4cm}p{3.4cm}}
\hline
\textbf{Współczynnik} & \textbf{Orientacja} & \textbf{FPS} & \textbf{Auto-detekcja} \\
\hline
Kadencja            & nie & tak & tak (z~metadanych FPS) \\
GCT                 & nie & tak & tak \\
Czas lotu           & nie & tak & tak \\
Stride length       & nie & tak & tak \\
Duty factor         & nie & tak & tak \\
Kąty stawów         & minimalnie & nie & nie (cichy artefakt) \\
Oscylacja pionowa   & tak (out-of-plane) & nie & częściowo \\
Foot strike         & tak (krytycznie) & nie & tak (próg $45^\circ$) \\
Symetria L/P        & tak (perspektywa) & nie & częściowo (SI~$>$~50\% jako przesłanka) \\
\hline
\end{tabular}
\end{table}

Z~tabeli wynika wyraźny podział. Metryki temporalne są odporne na orientację, lecz wrażliwe na FPS --- a~FPS jest znany z~metadanych pliku, więc problem da się wykryć w~pełni automatycznie. Metryki przestrzenne są odwrotnie: niewrażliwe na FPS, ale wrażliwe na perspektywę. Tu auto-detekcja jest trudniejsza, bo nie ma metadanych mówiących o~poprawności ujęcia --- trzeba wnioskować z~samych wartości. Najtrudniejszy przypadek to kąty stawów, których zniekształcenie perspektywiczne jest „ciche": wartość pozostaje w~pozornie sensownym zakresie, więc nie sposób jej odrzucić bez pomiaru referencyjnego.

\section{Automatyczna detekcja niewiarygodnych pomiarów}
\label{sec:auto-detekcja}

Trzecia hipoteza --- możliwość automatycznego wykrycia niewiarygodnych pomiarów bez ground truth --- znajduje potwierdzenie w~przypadkach, gdzie wartość metryki wykracza poza zakres fizjologicznie możliwy. Zaimplementowałem dwa takie mechanizmy.

\textbf{Flaga \texttt{low\_confidence} dla wzorca lądowania stopy.} W~funkcji obliczającej wzorzec lądowania oznaczam wynik jako niewiarygodny, gdy średni kąt stopy przekracza $45^\circ$ w~wartości bezwzględnej. Próg ten skalibrowałem na podstawie walidacji wizualnej ze~Studium~1: Janek (poprawne ujęcie) osiągał $|\text{kąt}| \le 12^\circ$, Adam (ujęcie graniczne) $12$--$33^\circ$, a~Film~10 (ujęcie pionowe) $82$--$160^\circ$. Próg $45^\circ$ oddziela „poprawne z~marginesem" od „ewidentny artefakt". Gdy flaga się zapala, silnik reguł zwraca ostrzeżenie o~niskiej wiarygodności i~wstrzymuje reguły interpretujące wzorzec lądowania --- nie ma sensu doradzać na podstawie wartości, która jest artefaktem perspektywy.

\textbf{Reguła łączona jakości predykcji.} Drugi mechanizm, opisany w~sekcji~\ref{sec:rules-combined}, wychwytuje sytuacje, w~których wysoki wskaźnik symetrii łączy się z~obniżoną pewnością predykcji --- sygnał, że asymetria L/P jest raczej artefaktem klasyfikatora niż cechą biegacza.

Oba mechanizmy działają bez pomiarów referencyjnych: opierają się wyłącznie na tym, czy wartość mieści się w~zakresie możliwym fizycznie i~czy różne wskaźniki jakości są ze~sobą spójne. To potwierdza trzecią hipotezę w~zakresie metryk, których zniekształcenie wyprowadza wartość poza zakres fizjologiczny. Ograniczeniem pozostają „ciche" artefakty kątów stawów, dla których taka detekcja nie jest możliwa.

\section{Walidacja na nagraniu własnym z~pomiarem referencyjnym}
\label{sec:walidacja-garmin}

Dotychczasowe oceny opierały się na porównaniu wyliczonych wartości z~zakresami z~literatury, ponieważ zbiór nie zawierał nagrań z~równoległym, niezależnym pomiarem. Aby uzupełnić ten brak, zarejestrowałem dodatkowe nagranie własnego biegu na bieżni mechanicznej (oznaczone w~danych jako \textit{Wojtek}, nagranie spoza zbiorów treningowego, walidacyjnego i~testowego), rejestrując jednocześnie dane z~zegarka sportowego Garmin. Zegarek zapisuje sekundowe wartości dynamiki biegu --- kadencję, czas kontaktu z~podłożem, długość kroku --- które posłużyły jako pomiar referencyjny (ground truth) dla wyników pipeline'u. Jest to jedyny w~całej pracy punkt odniesienia pochodzący z~zewnętrznego urządzenia, a~nie z~zakresów literaturowych.

Nagranie zarejestrowałem w~rozdzielczości $1920 \times 1080$ (poziomej) przy 29{,}97~FPS. Pełna aktywność obejmowała rozbieganie i~próby nagrania; do analizy wykorzystałem fragment biegu właściwego ($\approx$~35~s, kadencja ustabilizowana). Porównania dokonałem względem odpowiadającego mu okna danych Garmina --- średnia z~całej aktywności byłaby zaniżona przez wolne rozbieganie i~przerwy. Tabela~\ref{tab:garmin-validation} zestawia wyniki.

\begin{table}[htbp]
\centering
\caption{Porównanie pipeline'u z~pomiarem referencyjnym Garmin (okno biegu właściwego)}
\label{tab:garmin-validation}
\small
\begin{tabular}{lccc}
\hline
\textbf{Metryka} & \textbf{Pipeline} & \textbf{Garmin} & \textbf{Różnica} \\
\hline
Kadencja [spm]   & 162{,}5 & 163{,}3 & $-0{,}5\%$ \\
Czas cyklu [ms]  & 737     & $\approx$735 & $+0{,}3\%$ \\
GCT (śr. L/P) [ms] & 255   & 297     & $-14\%$ \\
\hline
\end{tabular}
\end{table}

Wynik ujawnia wyraźny kontrast, który stanowi empiryczne potwierdzenie analizy z~sekcji~\ref{sec:wplyw-na-metryki}. Kadencja i~czas cyklu zgadzają się z~Garminem niemal idealnie (poniżej 1\% różnicy), natomiast czas kontaktu z~podłożem jest systematycznie zaniżony o~około 14\%. Różnica ta wynika z~natury obu metryk: kadencja jest \textbf{licznikiem} kontaktów --- wystarczy wykryć, że kontakt nastąpił, niezależnie od tego, gdzie dokładnie wypada granica między fazą kontaktu a~lotem. Czas kontaktu jest natomiast \textbf{miarą czasu trwania} fazy STANCE, więc zależy bezpośrednio od położenia tej granicy. Pipeline systematycznie skraca fazę STANCE (zaniżając jednocześnie duty factor: 0{,}345 wobec wartości $\approx$0{,}40 wynikającej z~pomiaru Garmina), co wynika z~przyjętego w~auto-etykietowaniu parametru \texttt{flight\_fraction}~=~0{,}4 oraz z~różnicy definicyjnej (Garmin wyznacza kontakt akcelerometrem, od~uderzenia pięty do~oderwania palców). Sprawdzenie to dostarcza ilościowego dowodu na to, że metryki oparte na zliczaniu zdarzeń są znacznie odporniejsze niż metryki oparte na pomiarze czasu trwania faz.

Nagranie potwierdziło również wniosek ze~Studium~1. Mimo poziomej orientacji i~wysokiej rozdzielczości, kąt stopy w~momencie kontaktu wyniósł $-90{,}5^\circ$ (lewa) i~$-98{,}4^\circ$ (prawa) --- wartości fizycznie niemożliwe, automatycznie oznaczone flagą \texttt{low\_confidence}. Pokazuje to, że pozioma orientacja jest warunkiem koniecznym, ale niewystarczającym wiarygodności wzorca lądowania: kamera musi być również ustawiona prostopadle do toru biegu, czego w~tym nagraniu zabrakło. Współczynniki temporalne, niezależne od perspektywy, pozostały przy tym wiarygodne.

Warto odnotować dwie dodatkowe obserwacje. Po pierwsze, symetria lewo-prawo tego biegu była najlepsza w~całym materiale (maksymalny wskaźnik symetrii Robinsona 9{,}5\%, średni 3{,}1\%), co odpowiada balansowi kroków 49,5/50,5\% raportowanemu przez aplikację Garmin Connect. Po drugie, średnia pewność predykcji klasyfikatora wyniosła 0{,}855 --- najniższa spośród analizowanych biegaczy, choć wciąż powyżej progu ostrzeżenia 0{,}85.

\section{Podsumowanie}
\label{sec:podsumowanie-wrazliwosc}

Trzy studia przypadku potwierdziły trzy postawione hipotezy: współczynniki przestrzenne są wrażliwe na perspektywę kamery (Studium~1), temporalne na częstotliwość klatek (Studium~2), a~najbardziej zawodne przypadki --- wartości wykraczające poza zakres fizjologiczny --- da się wykryć automatycznie (sekcja~\ref{sec:auto-detekcja}). Walidacja na nagraniu własnym z~pomiarem referencyjnym Garmin (sekcja~\ref{sec:walidacja-garmin}) dodatkowo potwierdziła te wnioski niezależnym ground-truthem: kadencja zgodna w~granicach 1\%, czas kontaktu zaniżony o~14\% (potwierdzenie wrażliwości metryk czasu trwania), a~wzorzec lądowania automatycznie oznaczony jako niewiarygodny mimo poprawnej rozdzielczości i~orientacji.

Wkładem tej części pracy jest systematyczna analiza wrażliwości pipeline'u monocular 2D dla \textit{biegu}. Literatura badała podobne ograniczenia dla \textit{chodu} --- Stenum i~wsp.~\cite{stenum2021} opisali zależność długości kroku od pozycji w~kadrze, a~Menychtas i~wsp.~\cite{menychtas2023} błędy dla ruchów poza płaszczyzną obrazu --- ale dla biegu sportowego brakowało takiego zestawienia wraz z~działającym mechanizmem automatycznego ostrzegania użytkownika.

\chapter{Dyskusja i~wnioski}
\label{ch:dyskusja-wnioski}

\section{Odpowiedzi na pytania badawcze}
\label{sec:odpowiedz-pytania}

Pierwsze pytanie dotyczyło tego, jak zaprojektować pipeline wyliczający współczynniki biegu z~pojedynczego nagrania jedną kamerą i~jaką osiąga dokładność. Praca pokazała, że trójetapowy pipeline --- estymacja pozy przez MediaPipe, klasyfikacja faz siecią BiLSTM i~obliczenia geometryczne --- pozwala wyliczyć wszystkie dwanaście współczynników z~jednego nagrania (rozdział~\ref{ch:materialy-metody}). Dokładność oceniłem względem zakresów z~literatury, a~dla jednego nagrania także względem niezależnego pomiaru z~zegarka i opaski sportowej (rozdział~\ref{ch:wrazliwosc}). Kadencja zgodziła się z~pomiarem referencyjnym w~granicach 1\%, natomiast czas kontaktu z~podłożem był systematycznie zaniżony o~około 14\%. 

Drugie pytanie dotyczyło wyboru klasyfikatora faz i~jego wpływu na precyzję współczynników czasowych. Model sekwencyjny, uwzględniający kontekst sąsiednich klatek, przewyższył klasyfikator analizujący każdą klatkę niezależnie (rozdział~\ref{ch:klasyfikator}). Inżynieria cech poprawiła wyniki lasu losowego, ale większy zysk przyniosło dopiero wprowadzenie kontekstu czasowego w~sieci BiLSTM. Co istotne, dokładność klasyfikacji przekłada się nierównomiernie na precyzję poszczególnych metryk. Współczynniki oparte na zliczaniu zdarzeń, takie jak kadencja, okazały się odporne, podczas gdy współczynniki oparte na pomiarze czasu trwania faz, takie jak czas kontaktu z~podłożem, są wrażliwe na błędy granic faz. Wniosek ten wykazałem najpierw rachunkowo, a~następnie potwierdziłem niezależnym pomiarem referencyjnym.

Trzecie pytanie dotyczyło wrażliwości współczynników na warunki nagrania i~możliwości automatycznego wykrywania przypadków niewiarygodnych. Wykazałem, że współczynniki przestrzenne --- zwłaszcza wzorzec lądowania stopy --- są wrażliwe na perspektywę kamery, a~współczynniki czasowe na częstotliwość klatek (rozdział~\ref{ch:wrazliwosc}). Najbardziej zawodne przypadki, w~których wartość wykracza poza zakres fizjologicznie możliwy, udało się wykrywać automatycznie, bez pomiaru referencyjnego --- mechanizm ten zadziałał poprawnie nawet na nagraniu w~wysokiej rozdzielczości i~orientacji poziomej, gdzie problem perspektywy nie był oczywisty.

\section{Wkład pracy}
\label{sec:wklad-pracy}

Wkład pracy odpowiada lukom w~literaturze zidentyfikowanym w~rozdziale~\ref{ch:sota}. Po pierwsze, praca dotyczy biegu sportowego, a~nie chodu klinicznego, który dominuje w~dotychczasowych badaniach --- bieg ma odmienne fazy i~inne wartości referencyjne. Po drugie, zamiast proponować pojedynczy model, systematycznie porównałem cztery konfiguracje klasyfikatora na świadomie małym, dobrze opisanym zbiorze. Po trzecie, zbadałem wpływ korekty proporcji obrazu jako kroku preprocessingu dla nagrań o~nietypowej orientacji, czego wcześniejsze prace nie raportowały. Po czwarte, uporządkowałem wartości referencyjne pochodzące z literatury i zaimplementowałem je w postaci działającego silnika reguł rekomendacji. Każda reguła zawiera przypisane źródło oraz poziom ważności komunikatu.
 Po piąte, zaproponowałem i~wdrożyłem mechanizm automatycznego wykrywania niskiej wiarygodności pomiaru bez dostępu do danych referencyjnych.

Dodatkowym, niezaplanowanym pierwotnie wkładem jest walidacja względem zegarka sportowego. Jest to jedyny w~pracy punkt odniesienia pochodzący z~zewnętrznego urządzenia, a~nie z~zakresów literaturowych, i~to on dostarczył ilościowego dowodu na różną odporność metryk opartych na zliczaniu i~na pomiarze czasu trwania.

\section{Ograniczenia}
\label{sec:ograniczenia-pracy}

\begin{itemize}
    \item \textbf{Niewielki zbiór danych.} Piętnaście nagrań pochodzących od~kilkunastu biegaczy to bardzo niewielki zbiór danych. Ogranicza to możliwość uogólniania wyników na całą populację biegaczy.
    \item \textbf{Brak pełnego pomiaru referencyjnego 3D.} Nie dysponowałem systemem motion capture, więc nie mogę podać dokładności w~jednostkach fizycznych dla większości nagrań. 
\item \textbf{Systematyczne zaniżenie czasu kontaktu.} Porównanie z pomiarem referencyjnym pokazało, że pipeline zaniża czas kontaktu z podłożem o około 14\%. Prawdopodobną przyczyną jest sposób wyznaczania granic między fazą kontaktu i lotu oraz różnica między metodami pomiaru: Garmin korzysta z czujników ruchu, natomiast pipeline wyznacza fazy na podstawie obrazu.
\\
    \item \textbf{Ograniczenia obrazu dwuwymiarowego.} Współczynniki kątowe wyliczane są w~płaszczyźnie obrazu, więc ujęcie inne niż ściśle z boku zniekształca ich wartości.
    \item \textbf{Ogólny charakter wartości referencyjnych.} Progi reguł pochodzą z~literatury dla zdrowych biegaczy i~mogą nie odpowiadać indywidualnej sytuacji konkretnej osoby.
\item \textbf{Drobne różnice między uruchomieniami MediaPipe.} Przy ponownym przetwarzaniu tego samego nagrania MediaPipe może zwrócić minimalnie inne pozycje keypointów. Zwykle są to bardzo małe różnice, ale mogą wpływać na wynik, jeśli dana wartość znajduje się blisko progu reguły diagnostycznej.

\end{itemize}

\section{Warunki wiarygodnego działania systemu}
\label{sec:implikacje}

Z~analizy wynikają jasne wskazania, kiedy systemowi można ufać. Przy standardowym ujęciu z~boku (orientacja pozioma, kamera prostopadła do toru biegu) i~częstotliwości co najmniej 15~klatek na sekundę wiarygodne są zarówno współczynniki czasowe, jak i~wskaźniki symetrii. Przy ujęciu pionowym lub ukośnym wzorzec lądowania stopy staje się niewiarygodny, o~czym system ostrzega automatycznie. Przy niskiej częstotliwości klatek metryki czasowe tracą precyzję. System odpowiada na inną potrzebę niż laboratoryjne systemy referencyjne: udostępnia obiektywną analizę techniki biegu każdemu, kto dysponuje zwykłą kamerą --- bez kosztownego sprzętu, laboratorium i~obecności specjalisty. Jego realną siłą jest połączenie tej dostępności z~automatycznym sygnalizowaniem, którym wartościom można zaufać, a~które wymagają ostrożności --- cecha, której tradycyjne narzędzia pomiarowe nie oferują.

\section{Kierunki dalszych badań}
\label{sec:praca-przyszla}

Pracę można rozwijać w~kilku kierunkach.

\begin{itemize}
    \item \textbf{Większy i~bardziej zróżnicowany zbiór danych.} Rozmiar zbioru był głównym ograniczeniem pracy. Zebranie nagrań od~znacznie większej liczby biegaczy --- o~różnej technice, prędkości i~budowie ciała --- pozwoliłoby rzetelniej ocenić generalizację, ograniczyć przeuczenie klasyfikatora i~przejść na podział danych na poziomie biegaczy.
    \\
    \item \textbf{Dokładniejsze źródło keypointów.} MediaPipe wprowadza szum, który ogranicza precyzję współczynników, zwłaszcza kątowych. Warto zbadać dokładniejsze metody estymacji pozy, a~na etapie badawczym pozyskać keypointy z~systemu motion capture jako materiał o~wyższej jakości --- posłużyłyby one zarówno do zmierzenia, ile precyzji tracimy przez szum estymacji, jak i~jako lepszy punkt odniesienia przy dalszym rozwoju modelu.
    \item \textbf{Rekalibracja granic faz.} Wykryte w~porównaniu z~pomiarem Garmin systematyczne zaniżenie czasu kontaktu o~około 14\% wskazuje konkretny kierunek poprawy --- dostrojenie sposobu wyznaczania granic między fazą kontaktu a~fazą lotu, tak aby lepiej odpowiadały rzeczywistym momentom kontaktu i~oderwania stopy.
    \item \textbf{Przeliczanie wyników na jednostki fizyczne.} Wprowadzenie kalibracji skali obrazu (np.~na podstawie znanego wzrostu biegacza) pozwoliłoby raportować współczynniki takie jak oscylacja pionowa czy długość kroku w~centymetrach i~metrach zamiast w~wartościach znormalizowanych.
\end{itemize}

Pipeline można też uzupełnić o~automatyczne ostrzeganie użytkownika o~zbyt niskiej częstotliwości klatek lub niewłaściwej orientacji nagrania już na etapie wczytywania wideo --- co wynika bezpośrednio z~przeprowadzonej analizy wrażliwości.

\section{Wnioski końcowe}
\label{sec:wnioski-koncowe}

Praca wykazała, że z~pojedynczego nagrania biegu zwykłą kamerą można, z~użyciem uczenia maszynowego do klasyfikacji faz, wyliczyć komplet współczynników biomechanicznych i~wygenerować na ich podstawie rekomendacje oparte na literaturze. Kluczowym wynikiem jest rozróżnienie współczynników odpornych od~wrażliwych na warunki nagrania oraz wykazanie, że te najbardziej zawodne można rozpoznać automatycznie. System nie dorównuje dokładnością laboratoryjnym systemom referencyjnym, ale oferuje to, czego one nie dają --- dostępność i~uczciwą, automatyczną sygnalizację granic własnej wiarygodności.


% ============================================================================
% BIBLIOGRAFIA
% ============================================================================
% Template SGGW używa ręcznego \thebibliography (nie BibTeX/biblatex).
% Wpisy w chapters/bibliography.tex dla łatwiejszego zarządzania.

% ============================================================================
% BIBLIOGRAFIA
% ============================================================================
% Template SGGW używa ręcznego \thebibliography (nie BibTeX/biblatex).
% Wpisy uporządkowane tematycznie: biomechanika biegu / pose estimation / ML / metodologia.
%
% Status: 27 wpisów (sesja 2026-05-23). Zweryfikowane 3x (agenci) + ręcznie (user).
% Diaz 2019 usunięty — niepotwierdzony. OpenPose rok zmieniony na 2021 (tom 43, nr 1).
% Materiał źródłowy: docs/thesis-notes/2026-05-14-research-podobnych-prac.md
%                    docs/reference-values.md (cytaty biomechaniczne w rules.py)

\begin{thebibliography}{99}

% ============================================================================
% Pose estimation z monocular 2D wideo
% ============================================================================

\bibitem{bazarevsky2020blazepose}
V.\ Bazarevsky, I.\ Grishchenko, K.\ Raveendran, T.\ Zhu, F.\ Zhang, M.\ Grundmann,
\textit{BlazePose: On-device Real-time Body Pose Tracking},
CVPR Workshop on Computer Vision for Augmented and Virtual Reality, arXiv:2006.10204, 2020.

\bibitem{cao2021openpose}
Z.\ Cao, G.\ Hidalgo, T.\ Simon, S.E.\ Wei, Y.\ Sheikh,
\textit{OpenPose: Realtime Multi-Person 2D Pose Estimation Using Part Affinity Fields},
IEEE Transactions on Pattern Analysis and Machine Intelligence, 43(1): 172--186, 2021.
\url{https://doi.org/10.1109/TPAMI.2019.2929257}

\bibitem{lin2014coco}
T.Y.\ Lin et al.,
\textit{Microsoft COCO: Common Objects in Context},
European Conference on Computer Vision (ECCV), pp.\ 740--755, 2014.
\url{https://doi.org/10.1007/978-3-319-10602-1_48}

% ============================================================================
% Walidacja 2D pose vs 3D motion capture (gait/running)
% ============================================================================

\bibitem{stenum2021}
J.\ Stenum, C.\ Rossi, R.T.\ Roemmich,
\textit{Two-dimensional video-based analysis of human gait using pose estimation},
PLOS Computational Biology, 17(4): e1008935, 2021.
\url{https://doi.org/10.1371/journal.pcbi.1008935}

\bibitem{ali2024dmd}
M.M.\ Ali, M.M.\ Hassan, M.\ Zaki,
\textit{Human Pose Estimation for Clinical Analysis of Gait Pathologies},
Bioinformatics and Biology Insights, 2024.
\url{https://doi.org/10.1177/11779322241231108}

\bibitem{menychtas2023}
D.\ Menychtas, N.\ Petrou, I.\ Kansizoglou, E.\ Giannakou, A.\ Grekidis, A.\ Gasteratos,
V.\ Gourgoulis, E.\ Douda, I.\ Smilios, M.\ Michalopoulou, G.C.\ Sirakoulis, N.\ Aggeloulis,
\textit{Gait analysis comparison between manual marking, 2D pose estimation algorithms, and 3D marker-based system},
Frontiers in Rehabilitation Sciences, 4: 1238134, 2023.
\url{https://doi.org/10.3389/fresc.2023.1238134}

\bibitem{hgcnmlp2024}
R.\ Hu, Y.\ Diao, Y.\ Wang, G.\ Li, R.\ He, Y.\ Ning, N.\ Lou, G.\ Li, G.\ Zhao,
\textit{Effective evaluation of HGcnMLP method for markerless 3D pose estimation of musculoskeletal diseases patients based on smartphone monocular video},
Frontiers in Bioengineering and Biotechnology, 11: 1335251, 2024.
% Uwaga: volume 11 (2023), ale publikacja online 9 stycznia 2024
\url{https://doi.org/10.3389/fbioe.2023.1335251}

\bibitem{frontiers2025review}
S.\ Edriss, C.\ Romagnoli, L.\ Caprioli, V.\ Bonaiuto, E.\ Padua, G.\ Annino,
\textit{Commercial vision sensors and AI-based pose estimation frameworks for markerless motion analysis in sports and exercises: a mini review},
Frontiers in Physiology, 16: 1649330, 2025.
\url{https://doi.org/10.3389/fphys.2025.1649330}

\bibitem{roggio2024review}
F.\ Roggio, B.\ Trovato, M.\ Sortino, G.\ Musumeci,
\textit{A comprehensive analysis of the machine learning pose estimation models used in human movement and posture analyses: A narrative review},
Heliyon, 10(21): e39977, 2024.
\url{https://doi.org/10.1016/j.heliyon.2024.e39977}

\bibitem{vincent2024footstrike}
H.K.\ Vincent, K.\ Coffey, A.\ Villasuso, K.R.\ Vincent, S.\ Sharififar, L.\ Pezzullo, R.M.\ Nixon,
\textit{Accuracy of self-reported foot strike pattern detection among endurance runners},
Frontiers in Sports and Active Living, 6: 1491486, 2024.
\url{https://doi.org/10.3389/fspor.2024.1491486}

% ============================================================================
% Biomechanika biegu — reference values dla reguł rekomendacji
% ============================================================================

\bibitem{heiderscheit2011}
B.C.\ Heiderscheit, E.S.\ Chumanov, M.P.\ Michalski, C.M.\ Wille, M.B.\ Ryan,
\textit{Effects of step rate manipulation on joint mechanics during running},
Medicine \& Science in Sports \& Exercise, 43(2): 296--302, 2011.
\url{https://doi.org/10.1249/MSS.0b013e3181ebedf4}

\bibitem{novacheck1998}
T.F.\ Novacheck,
\textit{The biomechanics of running},
Gait \& Posture, 7(1): 77--95, 1998.
\url{https://doi.org/10.1016/S0966-6362(97)00038-6}

\bibitem{souza2016}
R.B.\ Souza,
\textit{An Evidence-Based Videotaped Running Biomechanics Analysis},
Physical Medicine and Rehabilitation Clinics of North America, 27(1): 217--236, 2016.
\url{https://doi.org/10.1016/j.pmr.2015.08.006}

% Diaz 2019 (vertical oscillation) — usunięty: nie udało się potwierdzić istnienia pracy.
% Jeśli znajdziesz pełne dane bibliograficzne, dodaj wpis z powrotem.

\bibitem{robinson1987}
R.O.\ Robinson, W.\ Herzog, B.M.\ Nigg,
\textit{Use of force platform variables to quantify the effects of chiropractic manipulation on gait symmetry},
Journal of Manipulative and Physiological Therapeutics, 10(4): 172--176, 1987.

\bibitem{daoud2012}
A.I.\ Daoud et al.,
\textit{Foot strike and injury rates in endurance runners: a retrospective study},
Medicine \& Science in Sports \& Exercise, 44(7): 1325--1334, 2012.
\url{https://doi.org/10.1249/MSS.0b013e3182465115}

% ============================================================================
% Uczenie maszynowe / klasyfikatory faz biegu
% ============================================================================

\bibitem{jeon2024bilstm}
H.\ Jeon, D.\ Lee,
\textit{Bi-Directional Long Short-Term Memory-Based Gait Phase Recognition Method Robust to Directional Variations in Subject's Gait Progression Using Wearable Inertial Sensor},
Sensors, 24(4): 1276, 2024.
\url{https://doi.org/10.3390/s24041276}

\bibitem{su2020dcnn}
B.\ Su, C.\ Smith, E.\ Gutierrez~Farewik,
\textit{Gait Phase Recognition Using Deep Convolutional Neural Network with Inertial Measurement Units},
Biosensors, 10(9): 109, 2020.
\url{https://doi.org/10.3390/bios10090109}

\bibitem{vu2020review}
H.T.T.\ Vu, D.\ Dong, H.-L.\ Cao, T.\ Verstraten, D.\ Lefeber, B.\ Vanderborght, J.\ Geeroms,
\textit{A Review of Gait Phase Detection Algorithms for Lower Limb Prostheses},
Sensors, 20(14): 3972, 2020.
\url{https://doi.org/10.3390/s20143972}

\bibitem{dimmick2023fatigue}
H.L.\ Dimmick, C.R.\ van~Rassel, M.J.\ MacInnis, R.\ Ferber,
\textit{Use of subject-specific models to detect fatigue-related changes in running biomechanics: a random forest approach},
Frontiers in Sports and Active Living, 5: 1283316, 2023.
\url{https://doi.org/10.3389/fspor.2023.1283316}

\bibitem{martinez2020rf}
J.\ Mart\'{i}nez-Gramage, J.P.\ Albiach, I.N.\ Molt\'{o}, J.J.\ Amer-Cuenca, V.H.\ Moreno, E.\ Segura-Ort\'{i},
\textit{A Random Forest Machine Learning Framework to Reduce Running Injuries in Young Triathletes},
Sensors, 20(21): 6388, 2020.
\url{https://doi.org/10.3390/s20216388}

\bibitem{young2023smartphone}
F.\ Young, R.\ Mason, R.\ Morris, S.\ Stuart, A.\ Godfrey,
\textit{Internet-of-Things-Enabled Markerless Running Gait Assessment from a Single Smartphone Camera},
Sensors, 23(2): 696, 2023.
\url{https://doi.org/10.3390/s23020696}

% ============================================================================
% Filtrowanie sygnałów biomechanicznych
% ============================================================================

\bibitem{crenna2021filtering}
F.\ Crenna, G.B.\ Rossi, M.\ Berardengo,
\textit{Filtering Biomechanical Signals in Movement Analysis},
Sensors, 21(13): 4580, 2021.
\url{https://doi.org/10.3390/s21134580}

\bibitem{savitzky1964}
A.\ Savitzky, M.J.E.\ Golay,
\textit{Smoothing and Differentiation of Data by Simplified Least Squares Procedures},
Analytical Chemistry, 36(8): 1627--1639, 1964.
\url{https://doi.org/10.1021/ac60214a047}

% ============================================================================
% Prewencja urazów biegowych / gait retraining
% ============================================================================

\bibitem{figueiredo2025cadence}
I.\ Figueiredo, M.\ Reis~e~Silva, J.E.\ Sousa,
\textit{The Influence of Running Cadence on Biomechanics and Injury Prevention: A Systematic Review},
Cureus, 17(8): e90322, 2025.
\url{https://doi.org/10.7759/cureus.90322}

% ============================================================================
% Uczenie maszynowe — podręczniki
% ============================================================================

\bibitem{goodfellow2016deep}
I.\ Goodfellow, Y.\ Bengio, A.\ Courville,
\textit{Deep Learning},
MIT Press, 2016.

\bibitem{hochreiter1997lstm}
S.\ Hochreiter, J.\ Schmidhuber,
\textit{Long Short-Term Memory},
Neural Computation, 9(8): 1735--1780, 1997.
\url{https://doi.org/10.1162/neco.1997.9.8.1735}

\bibitem{breiman2001rf}
L.\ Breiman,
\textit{Random Forests},
Machine Learning, 45(1): 5--32, 2001.
\url{https://doi.org/10.1023/A:1010933404324}

% ============================================================================
% Narzędzia / framework
% ============================================================================

\bibitem{pytorch}
A.\ Paszke et al.,
\textit{PyTorch: An Imperative Style, High-Performance Deep Learning Library},
Advances in Neural Information Processing Systems (NeurIPS), 2019.

\bibitem{scipy}
P.\ Virtanen et al.,
\textit{SciPy 1.0: fundamental algorithms for scientific computing in Python},
Nature Methods, 17: 261--272, 2020.
\url{https://doi.org/10.1038/s41592-019-0686-2}

\bibitem{sklearn}
F.\ Pedregosa et al.,
\textit{Scikit-learn: Machine Learning in Python},
Journal of Machine Learning Research, 12: 2825--2830, 2011.

\bibitem{opencv}
G.\ Bradski,
\textit{The OpenCV Library},
Dr. Dobb's Journal of Software Tools, 2000.

\end{thebibliography}


% ============================================================================
% PRZEDOSTATNIA STRONA (oświadczenie końcowe)
% ============================================================================

\beforelastpage

\end{document}
